\section{子拓扑斯与拓扑斯的粘合}\etiquette[9]{IV.sec9}\pageoriginale

\subsection{子拓扑斯与拓扑斯的嵌入}\etiquette[9.1]{IV.subsec9.1}

\begin{enonce*}{定义 9.1.1}\etiquette[9.1.1]{IV.def9.1.1}
\begin{enumerate}
\renewcommand{\theenumi}{\alph{enumi}}
\renewcommand{\labelenumi}{\theenumi)}
\item 若函子 $f_{\ast}:F\rightarrow E$ 全忠实
(\ie 伴随态射 $f^{\ast}f_{\ast}\rightarrow\id_F$ 是同构),
就称拓扑斯态射 $f:F\rightarrow E$ 是一个嵌入.

\item 拓扑斯 $E$ 的一个子拓扑斯,
是 $E$ 的一个本身为拓扑斯的严格全子范畴 $E'$,
并且其包含函子 $\alpha:E'\rightarrow E$ 形如 $i_{\ast}$,
其中 $i:E'\rightarrow E$ 是一个拓扑斯态射
(\ie (\hyperref[IV.subsubsec3.1.1]{3.1.1})
$\alpha$ 有一个左正合的左伴随 $i^{\ast}$).
此时态射 $i$ (它在相差唯一同构的意义下确定)
称为子拓扑斯 $E'$ 到 $E$ 中的包含态射.
\end{enumerate}
\end{enonce*}

\begin{enonce*}[remark]{注 9.1.2}\etiquette[9.1.2]{IV.rem9.1.2}
\begin{enumerate}
\renewcommand{\theenumi}{\alph{enumi}}
\renewcommand{\labelenumi}{\theenumi)}
\item 显然, 子拓扑斯的包含态射是一个嵌入,
且拓扑斯 $E$ 的一个子范畴 $E'$ 是子拓扑斯,
当且仅当它是某个嵌入拓扑斯态射 $f:F\rightarrow E$
所对应的正像函子 $f_{\ast}$ 的本质像.

\item 直接由定义可知, 拓扑斯态射 $f:F\rightarrow E$ 是嵌入,
当且仅当它同构于一个复合态射
$$
F\xrightarrow{g}E'\xrightarrow{i}E{\quoi},
$$
其中 $g$ 是一个拓扑斯等价,
$i$ 是 $E$ 的一个子拓扑斯 $E'$ 的包含态射.
后者由上述条件唯一确定, 即为 $f_{\ast}$ 的本质像,
上述分解也在相差唯一同构的意义下唯一.
当然, 在实际使用中, 应通过 $g$ 将 $F$ 与 $E$ 的子拓扑斯 $E'$ 等同
(比较 \hyperref[IV.subsubsec3.4.1]{3.4.1}).

\item 拓扑斯等价 (\hyperref[IV.subsec3.4]{3.4})
$u:E\xrightarrow{\approx}E_1$
按显然的方式定义 $E$ 的子拓扑斯所成集合与 $E_1$ 的子拓扑斯所成集合之间的一个双射.
事实上, 由于 $u_{\ast}$ 是范畴等价,
它在 $E$ 的严格全子范畴所成集合与 $E_1$ 的严格全子范畴所成集合之间
建立了\pageoriginale 一个双射
(与 $E'\subset E$ 相应的是 $u_{\ast}|E'$ 的本质像 $E'_1$),
而且直接可见, $E'$ 是 $E$ 的子拓扑斯,
当且仅当 $E'_1$ 是 $E_1$ 的子拓扑斯.
因此在研究子拓扑斯时, 可以归结到 $E$ 形如 $\widetilde{C}$ 的情形,
其中 $C$ 是一个小景.
在此情形中, \hyperrefext[II][II.theo5.5]{5.5} 表明:
$E$ 的一个严格全子范畴 $E'$ 是子拓扑斯,
当且仅当它的包含函子 $\alpha:E'\rightarrow E$
有一个左正合的左伴随 (这已经蕴含 $E'$ 是拓扑斯);
并且这些子拓扑斯与 $C$ 上比 $C$ 的给定拓扑 $T$ 更细的拓扑 $T'$ 一一对应:
与这样的 $T'$ 相应的是 $C^{\sim}$ 的严格全子范畴 $(C,T')^{\sim}$,
其对象是拓扑 $T'$ 的层.
这还表明, 对\emph{每个} $𝒰$-拓扑斯 $E$,
$E$ 的一个严格全子范畴 $E'$ 是子拓扑斯,
当且仅当它的包含函子 $\alpha:E'\rightarrow E$ 有一个左正合的左伴随;
\emph{$E$ 的子拓扑斯所成的有序集 $\mathfrak{E}(E)$ 是 $𝒰$-小的,
且 $\mathfrak{E}(E)$ 的每个子集都有上确界和下确界.}

\item 若 $E'$ 和 $E''$ 是拓扑斯 $E$ 的两个子拓扑斯,
则它们的交也是 $E$ 的子拓扑斯.
事实上, 由于这个交是严格全的, 可归结到 $E$ 形如 $C^{\sim}$ 的情形.
沿用 \sisi{b)}{c)} 的记号,
此时 $E'$ 和 $E''$ 对应于 $C$ 上两个比 $T$ 更细的拓扑 $T'$ 和 $T''$;
只须证明 \emph{$E'\cap E''$ 是拓扑 $T'''$ 的层范畴,
其中它是 $T'$ 和 $T''$ 的上确界}.
我们指出这个事实的证明
(本应作为推论放在 \hyperrefext[II][II.theo4.4]{4.4}
或 \hyperrefext[II][II.theo5.5]{5.5} 之后!).
设 $E'''$ 是 $T'''$ 的层范畴,
它显然包含在 $E'$ 和 $E''$ 中, 从而包含在 $E'\cap E''$ 中.
对 $E$ 的每个全子范畴 $D$, 设 $t(D)$ 是使 $D$ 的对象均为层的拓扑中最细的一个
\pageoriginale (\hyperrefext[II][II.prop2.2]{2.2}).
包含关系 $E'''\subset E'\cap E''\subset E',E''$ 显然蕴含
$t(E'),t(E'')\leq t(E'\cap E'')\leq t(E''')$.
另一方面,
(\sisi{\hyperrefext[II][II.theo4.4]{4.4} I}{\hyperrefext[II][II.cor4.4.4]{4.4.4}})
$t(E')=T'$, $t(E'')=T''$ 且 $t(E''')=T'''$,
所以上述不等式以及 $T'''$ 作为上确界的定义蕴含
$T'''=t(E'\cap E'')$, 从而 $E'\cap E''\subset E'''$, \cqfd\
更一般地, 这个论证表明:
若 $T'$ 是任意一个拓扑族 $(T_i)$ 的上确界,
则 $T'$ 的层范畴 $E'$ 是各个 $T_i$ 的层范畴 $E_i$ 的交.
特别地得到:
\end{enumerate}
\end{enonce*}

\begin{enonce*}[remark]{命题 9.1.3}\etiquette[9.1.3]{IV.prop9.1.3}
设 $E$ 是一个拓扑斯.
对 $E$ 的任意子拓扑斯族 $(E_i)_{i\in I}$,
交 $\bigcap_{i\in I}E_i$ 是 $E$ 的一个子拓扑斯.
\end{enonce*}

\begin{enonce*}{命题 9.1.4}\etiquette[9.1.4]{IV.prop9.1.4}
设 $i:E'\rightarrow E$ 是一个拓扑斯嵌入, $F$ 是一个拓扑斯.
函子
\begin{equation*}
f\longmapsto i\circ f:\Hhomtop(F,E')\longrightarrow
\Hhomtop(F,E)\tag{9.1.4.1}
\end{equation*}\etiquette[9.1.4.1]{IV.eq9.1.4.1}
全忠实, 且其本质像由满足 $g_{\ast}$ 的本质像包含在 $i_{\ast}$ 的本质像中的
拓扑斯态射 $g:F\rightarrow E$ 组成.
\end{enonce*}

事实上, 考虑显然的交换图
\[
\xymatrix@C=1.75cm@R=1cm@M=6pt{
\Hhomtop (F,E') \ar @{^{(}->}[d] \ar[r]& \Hhomtop (F,E) \ar @{^{(}->}[d]\\
\Hhom (F,E') \ar[r] &\Hhom (F,E) , }
\]
其中竖箭头是函子 $h\mapsto h_{\ast}$,
由 $\Hhomtop$ 的定义 (\sisi{IV}{}\hyperref[IV.subsec3.2]{3.2}), 它们全忠实.
第二个横箭头 $G\mapsto i_{\ast}G$ 也全忠实, 因为 $i_{\ast}$ 全忠实,
所以第一个横箭头也全忠实.
只须证明 \eqhyperref[IV.eq9.1.4.1]{9.1.4.1} 的本质像的刻画.
条件的必要性由公式 $(if)_{\ast}=i_{\ast}f_{\ast}$ 显然成立.
反过来, 假设 $g_{\ast}$ 的本质像包含在 $i_{\ast}$ 的本质像中,
\ie 可以写成 $g_{\ast}=i_{\ast}f_{\ast}$,
其中 $f_{\ast}:F\rightarrow E'$ 是一个函子.
须证明 $g$ 属于 \eqhyperref[IV.eq9.1.4.1]{9.1.4.1} 的本质像.
这就等于说 $f_{\ast}$ 有一个左正合的左伴随.
现在\pageoriginale 显然 $f_{\ast}$ 有左伴随 $g^{\ast}i_{\ast}$,
而这个函子左正合, 因为它是左正合函子 $g^{\ast}$ 和 $i_{\ast}$ 的复合, \cqfd

关于 \hyperref[IV.prop9.1.4]{9.1.4} 的逆命题,
\cf 习题 \hyperref[IV.exer9.1.6]{9.1.6}.

取 $F$ 为点拓扑斯 (\sisi{IV~}{}\hyperref[IV.subsec2.2]{2.2}),
由 \hyperref[IV.prop9.1.4]{9.1.4} 尤其得到:

\begin{enonce*}{推论 9.1.5}\etiquette[9.1.5]{IV.cor9.1.5}
若 $f:E'\rightarrow E$ 是一个嵌入拓扑斯态射,
则相应函子
\begin{equation*}
\Point(f):\Ppoint(E')\longrightarrow \Ppoint(E)\tag{9.1.5.1}
\end{equation*}\etiquette[9.1.5.1]{IV.eq9.1.5.1}
全忠实.
\end{enonce*}

\begin{enonce*}[remark]{习题 9.1.6}\etiquette[9.1.6]{IV.exer9.1.6}
(\emph{拓扑斯的嵌入与 2-单态射}).

设 $i:E'\rightarrow E$ 是一个拓扑斯态射.

\begin{enumerate}
\renewcommand{\theenumi}{\alph{enumi}}
\renewcommand{\labelenumi}{\theenumi)}
\item 若 $F$ 和 $G$ 是两个在 $E'$ \emph{之上}的拓扑斯,
定义 (沿用 \hyperref[IV.exer5.14]{5.14} a) 的记号) 一个典范函子
\begin{equation*}
f\longmapsto f\ast i:\Hhomtop_{E'}(F,G)\longrightarrow
\Hhomtop_E(F,G){\quoi}.\tag{9.1.6.1}
\end{equation*}\etiquette[9.1.6.1]{IV.eq9.1.6.1}

\item 假设 $i$ 作为 $𝒱\text{-}𝒰$-拓扑斯的 2-范畴
(\hyperref[IV.subsubsec3.3.2]{3.3.2}) 的 1-箭头,
是一个 $2$-单态射,
\ie 对每个属于 $𝒱$ 的 $𝒰$-拓扑斯 $F$,
函子 \hyperref[IV.eq9.1.4.1]{(9.1.4.1)} 全忠实. % add a ref
证明: 对每对 $𝒰$-拓扑斯 $F,G$ (不一定在 $𝒱$ 中),
函子 \hyperref[IV.eq9.1.6.1]{(9.1.6.1)} 全忠实. %% add a ref

\item 证明 $i$ 是嵌入, 当且仅当 $i$ 是 $2$-单态射.
(对于充分性, 将 b) 应用于诱导拓扑斯 $E'_{/X'}$ 和 $E'_{/Y'}$,
并利用 \hyperref[IV.exer5.14]{5.14} c).)

\item 设 $g:F\rightarrow E$ 是一个拓扑斯态射,
并假设 $2$-纤维积 (\hyperref[IV.prop5.11]{5.11})
\sisi{$F'=\fprod{F^2}{E'}{E}$}{$F'=F\times_E^2E'$}
存在 (P. DELIGNE 已证明这个条件总是满足的
\nde{另见 [16], 命题 4.47.}).
设 $j:F'\rightarrow F$ 是投射态射.
证明: 若 $i$ 是嵌入, 则 $j$ 也是.
(利用 c), 并注意 2-单态射这一概念对 2-基变换的稳定性是形式的.)
当 $E'$ 是 $E$ 的子拓扑斯,
且 $i:E'\rightarrow E$ 是典范包含态射时,
$F'$ 在 $j_{\ast}$ 下的本质像
(它是 $F$ 的一个子拓扑斯, 与 2-纤维积构造中\pageoriginale 的不确定性无关)
称为\emph{$F$ 中作为 $E$ 的子拓扑斯 $E'$ 在 $g:F\rightarrow E$ 下的逆像的子拓扑斯}.

\item 沿用 d) 末尾的记号,
取 $F$ 中作为 $E$ 的子拓扑斯 $E'$ 的逆像的子拓扑斯作为 2-纤维积.
设 $X$ 是 $F$ 的一个对象.
证明: $X$ 属于 $F'$ 的必要条件是满足以下条件:
对 $F$ 的每个对象 $U$,
以 $g_U:F_{/U}\rightarrow E$ 表示 $g$ 诱导的拓扑斯态射,
都有 $g_{U\ast}(X|U)\in\Ob E'$, 其中
$X|U=(X\times U\xrightarrow{\pr_2}U)$ 是“$X$ 在 $U$ 上的限制”.
(注意: 若 $i_U:F_{/U}\rightarrow F$ 是典范包含,
则 $g_U=g\circ i_U$,
从而 $g_{U\ast}(X|U)=g_{\ast}(i_{U\ast}(X|U))=g_{\ast}(\Hhom(U,X))$,
其中 $\Hhom(U,X)$ 是下文 \hyperref[IV.prop10.1]{10.1} 中定义的对象.)

\noindent \emph{问题.}
反过来, 关于 $X$ 的上述条件是否足以保证 $X$ 属于 $F'$?
这等于问, 由这些对象 $X$ 组成的 $F$ 的严格全子范畴,
是否为 $F$ 的一个\emph{子拓扑斯}.

\item 设 $X$ 是 $E$ 的一个对象, $X'=i^{\ast}(X)$,
$i_{/X}:E'_{/X'}\rightarrow E_{/X}$ 是 $i$ 诱导的拓扑斯态射
(\eqhyperref[IV.eq5.10.1]{5.10.1}).
证明: 若 $i$ 是嵌入, 则 $i_{/X}$ 也是.
(利用 d) 和 \hyperref[IV.prop5.11]{5.11}.)
反过来, 若 $(X_{\alpha})_{\alpha}$ 是一个覆盖终对象的对象族,
且对每个 $\alpha$, $i_{/X_{\alpha}}$ 都是嵌入, 则 $i$ 也是嵌入.
\end{enumerate}
\end{enonce*}

% END INCLUDED SGA4_English_current_buildable_source_20260906/source/04/drafts/SGA4_IV_LINES_5254_5490_English.texfrag


% BEGIN INCLUDED SGA4_English_current_buildable_source_20260906/source/04/drafts/SGA4_IV_LINES_5491_5778_English.texfrag
\begin{enonce*}[remark]{习题 9.1.7}\etiquette[9.1.7]{IV.exer9.1.7}
(\emph{子拓扑斯与箭头的乘法系}).

\begin{enumerate}
\renewcommand{\theenumi}{\alph{enumi}}
\renewcommand{\labelenumi}{\theenumi)}
\item 设 $i:E'\rightarrow E$ 是一个拓扑斯态射,
$S$ 是 $E$ 中使 $i^{\ast}(u)$ 为同构的箭头 $u$ 所成的集合,
并考虑由 $i^{\ast}$ 诱导的典范函子
(\cf [2, 第 I 章] 或 \sisi{VI}{\hyperrefext[VI][VI.prop6.2]{6.2} 的证明}):
\begin{equation*}
\rho:E[S^{-1}]\longrightarrow E'{\quoi}.\tag{9.1.7.1}
\end{equation*}\etiquette[9.1.7.1]{IV.eq9.1.7.1}

证明 $i$ 是嵌入, 当且仅当上述函子 $\rho$ 是范畴等价;
并且此时 $S$ 既容许左分式演算, 又容许右分式演算.
(利用 [2, 第 I 章, 1.3].)
此时, $i_{\ast}$ 的本质像可由 $S$ 恢复:
它由满足以下条件的 $X\in \Ob \sisi{\not{E}}{(E)}$ 组成:
对 $S$ 中的每个\pageoriginale $u:Y\rightarrow Y'$,
映射 $\Hom(u,X):\Hom(Y',X)\rightarrow\Hom(Y,X)$ 都是双射.

\item 此后假设 i 是嵌入. 证明对每个拓扑斯 $F$, 函子
$$
f^{\ast}\longmapsto f^{\ast}i^{\ast}:\Hhom(E',F)\longrightarrow
\Hhom(E,F)
$$
诱导如下两个范畴之间的一个等价:
来自拓扑斯态射 $f:F\rightarrow E'$ 的 $f^{\ast}$
(\ie 左正合并有右伴随的函子) 所成的范畴,
以及来自拓扑斯态射 $g:F\rightarrow E$ 的函子 $g^{\ast}:E\rightarrow F$ 所成的范畴,
其中 $g^{\ast}$ 把 $S$ 的箭头映为同构.

\item 设 $S$ 是 $Fl$ $E$ 的一个子集, 其中 $E$ 是一个 $𝒰$-拓扑斯.
证明 $S$ 由 a) 的过程与 $E$ 的一个子拓扑斯 $E'$ 相应,
当且仅当 $S$ 满足以下条件:
\begin{itemize}
\item[ST 1)] $S$ 包含可逆箭头, 并且对复合和基变换稳定.

\item[ST 2)] 若复合 vu 属于 $S$, 则 $u\in S$ 当且仅当 $v\in S$.

\item[ST 3)] 若 $u:X\rightarrow Y$ 满足:
存在一个覆盖族 $Y_i\rightarrow Y$, $i\in I$,
使对每个 $i\in I$,
$u_i:X_i=\fprod{X}{Y_i}{Y}\to Y_{i}$ 都属于 $S$,
则 $u\in S$.
\end{itemize}

证明这给出了 $E$ 的子拓扑斯 $E'$ 所成的集合,
与满足条件 ST 1), ST 2), ST 3) 的 $Fl$ $E$ 的子集 $S$ 所成的集合之间的一一对应.
(给定 $S$, 配上一个比 $E$ 的典范拓扑 $T$ 更细的 $𝒰$-拓扑 $T'$,
其覆盖族 $(X_i\rightarrow X)_{i\in I}$
为那些使 $X_i$ 的像的包含态射 $X'\hookrightarrow X$ 满足 $\in S$ 的族.
证明 $T'$ 的层范畴等价于 $E$ 的一个子拓扑斯 $E'$,
并利用 \hyperrefext[II][II.theo4.4]{4.4} 证明它恢复了 $S$.)

\item 证明 $S$ 由 $S$ 中的\emph{单态射}组成的子集 $S_0$ 确定.
证明映射 $S\mapsto S_0$ 在如下两个集合之间建立一个双射:
$Fl$ $E$ 中满足 c) 的条件 ST 1) 至 ST 3) 的子集 $S$ 所成的集合,
以及 $Fl$ $E$ 中由单态射\pageoriginale 组成并满足同样的条件 ST 1) 至 ST 3)
的子集所成的集合.

\item 设 $(E'_i)_{i\in I}$ 是 $E$ 的一个子拓扑斯族,
$(S_i)_{i\in I}$ 是相应的 $Fl$ $E$ 的子集族.
证明 $\bigcap_{i\in I} S_i$ 满足 c) 的条件 ST 1) 至 ST 3),
从而对应于 $E$ 的一个子拓扑斯 $E'$;
并且后者是\emph{由 $E'_i$ 生成的子拓扑斯},
\ie 子拓扑斯 $E'_i$ 的\emph{上确界}
(\cf \hyperref[IV.rem9.1.2]{9.1.2} c)).

\item 设 $A$ 是 $Fl$ $E$ 的一个子集.
证明: 在包含 $A$ 且满足条件 ST 1) 至 ST 3) 的 $Fl$ $E$ 的子集中,
存在最小的一个 $S$,
且 $S$ 对应于 $E$ 的如下最大的子拓扑斯 $E'$:
$Fl$ $E$ 的相应乘法系包含 $A$.
证明 $S$ 等于 $Fl$ $E$ 的子集 $A_n$ $(n\geq0)$ 的并,
它们递归构造如下:
$A_0=A$, 而 $A_{n+1}$ 由以下四种类型 (i) 至 (iv) 的箭头 $u:X\rightarrow Y$ 组成:
\begin{enumerate}
\renewcommand{\theenumii}{\roman{enumii}}
\renewcommand{\labelenumii}{(\theenumii)}
\item $u$ 可逆, 或者是 $A_n$ 的两个箭头的复合,
或者由 $A_n$ 的一个箭头经基变换得到.

\item $u$ 是某个复合箭头的一个因子,
而另一个因子以及复合都属于 $A_n$.

\item $u$ 是一个箭头小族 $\in A_n$ 的余积.

\item 存在一个满态射 $Y'\rightarrow Y$,
使由 $u$ 经基变换所得的
$u'=X'=\fprod{X}{Y'}{Y}\rightarrow Y'$ 属于 $A_n$.
\end{enumerate}

\item 沿用 e) 的记号, 设 $A$ 是 $S_i$ 的并,
$S$ 是从 $A$ 经 f) 的过程所得的 $Fl$ $E$ 的子集,
$E'$ 是 $E$ 的相应子拓扑斯.
证明 $E'$ 是子拓扑斯 $E'_i$ 的交 (= 下确界).
\end{enumerate}
\end{enonce*}



\begin{enonce*}[remark]{习题 9.1.7.2}\etiquette[9.1.7.2]{IV.exer9.1.7.2}
(\emph{拓扑斯态射的典范分解}).
\begin{enumerate}
\renewcommand{\theenumi}{\alph{enumi}}
\renewcommand{\labelenumi}{\theenumi)}
\item 设 $f:F\longrightarrow E$ 是一个拓扑斯态射.
证明在相差同构的意义下, 可以找到 $f$ 的一个分解
\begin{equation*}
F\xrightarrow{f'}E'\xrightarrow{i} E,\tag{9.1.7.3}
\end{equation*}\etiquette[9.1.7.3]{IV.eq9.1.7.3}
其中 $f'$ \emph{保守}
(\ie (\hyperref[IV.subsubsec6.4.0]{6.4.0})
${f'}^{\ast}$ 保守, 或等价地忠实)\pageoriginale,
$i$ 是嵌入;
并且这个分解在相差等价的意义下唯一
(\cf \hyperref[IV.prop9.1.4]{9.1.4}).
为此, 引入 $E$ 中使 $f^{\ast}(u)$ 为同构的箭头 $u$ 所成的集合 $S_f\subset Fl$ $E$,
证明 $S_f$ 满足 \hyperref[IV.exer9.1.7]{9.1.7} c) 的条件,
取 $E'$ 为 $E$ 的相应子拓扑斯,
再由 \hyperref[IV.eq9.1.7.1]{(9.1.7.1)} 定义 ${f'}^{\ast}$;
根据 \hyperref[IV.exer9.1.7]{9.1.7} b), 它对应于一个拓扑斯态射 $f'$.

\item 由嵌入 $i$ 定义的 $E$ 的子拓扑斯,
称为\emph{$E$ 中的拓扑斯态射 $f$ 的像子拓扑斯}.
证明它是 $E$ 的如下子拓扑斯 $E_1$ 中最小的一个:
$f$ 在相差同构的意义下经过 $E_1$ 分解;
\ie (\hyperref[IV.prop9.1.4]{9.1.4}),
是包含 $f_{\ast}$ 的像 (或等价地, $f_{\ast}$ 的本质像) 的 $E$ 的子拓扑斯中最小的一个.
$f$ 是嵌入的充要条件是 $f$ 诱导 $F$ 与其像之间的一个等价 $f'$;
$f$ 保守的充要条件是其像为整个 $E$.

\item 假设 $f=\Top(f_0)$ 对应于清醒拓扑空间的一个连续映射
$f_0:Y\rightarrow X$ (\hyperref[IV.subsec4.1]{4.1}).
考虑 $f_0$ 的典范分解
$$
Y\xrightarrow{{f'}_0}X'=f_0(Y)\xrightarrow{i_0}X{\quoi},
$$
其中 $i_0$ 是包含.
证明此时 $f$ 的典范分解可与相应分解
\[
\xymatrix@C=2cm@M=.2cm{ \Top(Y) \ar[r]^{\Top(f'_0)}& \Top(X')
\ar@{^{(}->}[r]^{\Top(i_0)}& \Top(X). }
\]
等同. 设 $Z$ 是 $X$ 中包含 $X'=f_0(Y)$ 的最小清醒子空间,
\ie $X$ 中满足 $\{\bar{x}\}\cap X'$ 在 $\{\bar{x}\}$ 中稠密的点 $x$ 所成的集合.
(注意, 若 $X$ 的点都是闭的,
或 $X'$ 是局部 Noether 空间 $X$ 的一个\emph{可构造}子集, 则 $Z=X'$.)
则 $f$ 保守,
\ie $E=\Top(X)$ 中作为 $F=\Top(Y)$ 在 $f=\Top(f_0)$ 下的像的子拓扑斯等于 $E$ 本身,
当且仅当 $Z=X$,
\ie 当且仅当 $f_0(Y)$ 在 $X$ 中\emph{很稠密} (EGA IV 10.13).
这就明确了在何种意义下,
可以合理地把保守拓扑斯态射的概念看作\emph{满射}连续映射概念的自然推广.

\item 设\pageoriginale
\[
\xymatrix{
Y \ar[d]& Y'\ar[l] \ar[d]\\
X & X'\ar[l] }
\]
是拓扑空间的一个 Cartesian 图.
若 $Y$ 和 $X'$ 都是离散空间, 因而 $Y'$ 也离散,
则相应的拓扑斯图是一个 2-纤维积图
(\hyperref[IV.prop5.11]{5.11}).
(注意, 若不再假设 $Y$ 和 $X'$ 离散,
取 $Y$ 和 $X'$ 为 Hausdorff 空间 $X$ 的子空间即可有反例,
\cf \hyperref[IV.exer9.1.10]{9.1.10} c).)
由此给出一个拓扑斯 2-纤维积图的例子
\[
\xymatrix{
F \ar[d]_{f}& F'\ar[l] \ar[d]^{f'}\\
E & E'\ar[l] }
\]
使 $f$ 保守而 $f'$ 不保守;
更确切地, 使 $E'=$ 点拓扑斯, $F'=$ 空拓扑斯
(\hyperref[IV.subsec2.2]{2.2}).
(取 $Y$ 使 $Y$ 在 $X$ 中的像很稠密且 $\neq X$,
并取 $X'$ 为一个不属于此像的点.)

\item 设 $(E_i)_{i\in I}$ 是拓扑斯 $E$ 的一个子拓扑斯族.
证明 $E$ 是 $E_i$ 的上确界,
当且仅当嵌入态射族 $E_i\rightarrow E$ 保守
(\hyperref[IV.subsubsec6.4.0]{6.4.0}).
\end{enumerate}
\end{enonce*}

\begin{enonce*}[remark]{习题 9.1.8}\etiquette[9.1.8]{IV.exer9.1.8}
(\emph{子拓扑斯与 $E$ 的点. 拓扑空间的情形}).

设 $E$ 是一个拓扑斯.
\begin{enumerate}
\renewcommand{\theenumi}{\alph{enumi}}
\renewcommand{\labelenumi}{\theenumi)}
\item 设 $Z$ 是 $\Ppoint(E)$ 的一个全子范畴
(或等价地, 是 $\Ob\Ppoint(E)$ 的一个子集).
证明 $E$ 中被与 $z\in\Ob Z$ 相应的纤维函子映为双射的箭头所成的集合 $S(Z)$,
满足 \hyperref[IV.exer9.1.7]{9.1.7} c) 的条件,
从而定义 $E$ 的一个子拓扑斯 $E_Z$.
证明此子拓扑斯有足够多的点,
且 $Z$ 等价于 $\Ppoint(E_Z)$ 的一个全子范畴.
证明 $E$ 的每个有足够多点的子拓扑斯 $E'$ 都可以通过上述过程得到.
(取 $Z$ 为函子 \hyperref[IV.eq9.1.5.1]{(9.1.5.1)} 的本质像.)
证明由此得到\pageoriginale 一个双射,
一端是 $E$ 的有足够多点的子拓扑斯 $E'$ 所成的集合,
另一端是 $\Ppoint(E)$ 的如下严格全子范畴 $Z$ 所成集合的某个子集:
它们对直因子
(\sisi{\hyperrefext[I][I.sec10]{10}}{\hyperrefext[I][I.sec10.6]{10.6}})
及小的滤过投射极限稳定 (\cf \hyperref[IV.subsec7.3]{7.3}).
(问题: $\Ppoint(E)$ 的哪些子范畴以这种方式出现?
根据 \hyperref[IV.rem9.1.2]{9.1.2} c) 中的对应字典,
这本质上是 \hyperref[IV.prop6.5.3]{6.5.3} 中已经提出的问题的另一种形式.
注意, 除了刚指出的必要条件之外,
还存在关于 $\Point(E)$ 的拓扑性质的更微妙条件,
例如清醒性; \cf e).)

\item 设 $(Z_i)_{i\in I}$ 是 $\Ppoint(E)$ 的一个全子范畴族,
$Z$ 是这些 $Z_i$ 的并所成的全子范畴.
证明 a) 中定义的相应子拓扑斯 $E_Z$,
就是 $E$ 中由 $E_{Z_i}$ 生成的子拓扑斯.
(利用 \hyperref[IV.exer9.1.7]{9.1.7} e).)
由此得出: 由一族有足够多点的子拓扑斯生成的 $E$ 的子拓扑斯也有足够多的点.

\item 设 $f:E'\rightarrow E$ 是一个拓扑斯嵌入.
证明映射 $U\mapsto f^{\ast}(U)$
\begin{equation*}
\Ouv(E)\longrightarrow \Ouv(E')\tag{9.1.8.1}
\end{equation*}\etiquette[9.1.8.1]{IV.eq9.1.8.1}
是满射, 并得出: $f$ 在点的同构类上诱导的映射
\begin{equation*}
\Point(E')\longrightarrow \Point(E)\tag{9.1.8.2}
\end{equation*}\etiquette[9.1.8.2]{IV.eq9.1.8.2}
是从 $\Point(E')$ 到 $\Point(E)$ 的某个子空间的同胚,
这里取 \hyperref[IV.exer7.8]{7.8} a) 中定义的拓扑.
(对于单射性, 利用 \hyperref[IV.cor9.1.5]{9.1.5}.)

\item 设 $f:X'\rightarrow X$ 是一个连续映射,
给出一个拓扑斯态射 (\hyperref[IV.subsec4.1]{4.1})
$$
\Top(f):\Top(X')\longrightarrow \Top(X){\quoi}.
$$
证明后者是嵌入,
当且仅当 $X'$ 的拓扑是 $X$ 的拓扑在 $f$ 下的逆像
(\ie 当 $X'$ 是 Kolmogorov 空间, 例如清醒空间时,
当且仅当 $f$ 诱导 $X'$ 到 $X$ 的某个子空间的同胚).
(对于必要性,\pageoriginale 利用 c);
对于充分性, 利用 \hyperref[IV.exer9.1.7.2]{9.1.7.2} c).)

\item 设 $X$ 是一个清醒拓扑空间, $E=\Top(X)$.
利用 a) 和 d), 在如下两个有序集之间建立一个同构 $X'\mapsto E_{X'}$:
$X$ 的清醒子空间 $X'$ 所成的集合
(\ie 当 $X$ 为 Hausdorff 空间时, $X$ 的\emph{所有}子集所成的集合),
以及 $E$ 的有足够多点的子拓扑斯所成的集合.
证明 $E_{X'}$ 等价于 $\Top(X')$.
若 $(X'_i)_{i\in I}$ 是 $X$ 的一个清醒子空间族,
证明存在 $X$ 的最小清醒子空间 $X'$ 包含全部 $X'_i$,
当 I 有限或 $X$ 为 Hausdorff 空间时, 它等于 $X'_i$ 的并;
并且 $E_{X'}$ 是 $E$ 中由 $E_{X'_i}$ 生成的子拓扑斯.
(利用 b).)

\item 给出一个形如 $\Top(X)$ 的拓扑斯 $E$ 及其子拓扑斯 $E'$ 的例子,
使后者没有足够多的点,
甚至非空 (\hyperref[IV.subsec2.2]{2.2}) 却没有点.
(取形如 $\widehat{U}$ 的 $E$ 及 $E'=U^{\sim}$,
其中 $U$ 是 \hyperref[IV.subsec7.4]{7.4} 中定义的有序集.
更好的办法是取 \sisi{9.1. 10. b)}{\hyperref[IV.exer9.1.10]{9.1.10} b)} 中的例子,
它表明可以取 $X$ 为任意一个没有孤立点的非空紧空间.)
\end{enumerate}
\end{enonce*}

% END INCLUDED SGA4_English_current_buildable_source_20260906/source/04/drafts/SGA4_IV_LINES_5491_5778_English.texfrag


% BEGIN INCLUDED SGA4_English_current_buildable_source_20260906/source/04/drafts/SGA4_IV_LINES_5779_6107_English.texfrag
\begin{enonce*}[remark]{习题 9.1.9}\etiquette[9.1.9]{IV.exer9.1.9}
(\emph{由有序集定义的拓扑斯的情形}).

\begin{enumerate}
\renewcommand{\theenumi}{\alph{enumi}}
\renewcommand{\labelenumi}{\theenumi)}
\item 使用 \hyperref[IV.exer7.8]{7.8} b) 和 c) 的方法,
它在如下两类对象之间给出对应字典:
一方面是由其开子对象生成的拓扑斯及这类拓扑斯的态射;
另一方面是容许任意 $\Sup$ 和有限下确界,
并且两个元素的下确界对任意 $\Sup$ 分配的有序集;
%% standardize notation sup/Sup/\Sup etc.
这类有序集之间的态射是与有限下确界及任意 $\Sup$ 交换的递增映射.

\item 设 $E$ 是一个拓扑斯, $f:E'\rightarrow E$ 是一个拓扑斯嵌入.
证明对 $E$ 中的每个生成族 $(X_i)_{i\in I}$,
族 $(f^{\ast}(X_i))_{i \in I}$ 也是生成族.
由此得出: 若 $E$ 由其开子对象生成, 则 $E'$ 也如此.
(利用 \hyperref[IV.exer9.1.8]{9.1.8} c).)

\item 设 $f:E'\rightarrow E$ 是由其开子对象生成的拓扑斯之间的一个态射,
对应于相应有序集之间的态射 $g:\Oo\rightarrow\Oo'$.
证明 $f$ 是嵌入, 当且仅当 $g$ 满射,
\emph{并且} $g$\pageoriginale 在序关系的图上也满射,
\ie $\Oo'$ 上的序是 $\Oo$ 上的序的商.

\item 设 $E$ 是 a) 中形如 $\Oo^{\sim}$ 的拓扑斯.
由 b) 和 c) 得出一个双射对应,
一端是 $E$ 的子拓扑斯所成的集合,
另一端是 $\Oo$ 上具有以下性质的等价关系 $R$ 所成的集合;
这些性质类似于 \hyperref[IV.exer9.1.7]{9.1.7} c) 的条件 ST 1) 至 ST 3):
\begin{itemize}
\item[Q0 1)] 关系 $R$ 与有限下确界相容;
换言之, 若 $U$ 和 $U'$ 在 $R$ 下等价,
则对每个 $V\in\Oo$, $U\cap V$ 和 $U'\cap V$ 也等价.

\item[Q0 2)] 关系 $R$ 与任意 $\Sup$ 相容;
换言之, 若对每个 $i\in I$, $U_i$ 与 $U'_i$ 等价,
则 $\Sup_i U_i$ 与 $\Sup_i U'_i$ 等价.
\end{itemize}

证明: 若 $E$ 的子拓扑斯 $E'_i$ 对应于等价关系 $R_i$,
则由 $E'_i$ 生成的子拓扑斯对应于关系 $R_i$ 的交.
\end{enumerate}
\end{enonce*}

\begin{enonce*}[remark]{习题 9.1.10}\etiquette[9.1.10]{IV.exer9.1.10}
(\emph{子空间的交; 新的无点拓扑斯; 子拓扑斯的补不存在}).

设 $X$ 是一个拓扑空间,
$E=\Top(X)=\Ouv(X)^{\sim}$ (\hyperref[IV.subsec2.1]{2.1}).

\begin{enumerate}
\renewcommand{\theenumi}{\alph{enumi}}
\renewcommand{\labelenumi}{\theenumi)}
\item 设 $R$ 为开集集合上的如下关系:
$\Ouv(X):(U,U')\in R\Longleftrightarrow U\cap U'$ 在 $U$ 和 $U'$ 中都稠密.
证明 $R$ 是一个满足 \hyperref[IV.exer9.1.9]{9.1.9} d)
中的条件 Q0 1) 和 Q0 2) 的等价关系,
因而定义 $E$ 的一个子拓扑斯 $E'$.
证明对一点 $x\in X$, 相应的点属于 $\Ppoint(E')$ 的本质像,
当且仅当 $x$ \emph{厚},
\ie 属于 $\bar{x}$ 的内部的闭包.
证明 $E'$ 是空拓扑斯 (\hyperref[IV.subsec2.2]{2.2}), 当且仅当 $X$ 为空.

\item 若 $X$ 清醒,
证明范畴 $\Ppoint(E')$ 等价于 $X$ 中由 $X$ 的厚点组成的有序子集
(取特殊化关系) 所定义的范畴.
证明若 $X$ 的点都是闭的 (\resp 若 $X$ 为 Noether 空间),
则 $X$ 的厚点是 $X$ 的孤立点,
\ie 使 $\{x\}$ 为开的点 (\resp 是 $X$ 的极大点).
由此得出: 若 $X$ 是\pageoriginale 一个没有孤立点的非空 Hausdorff 拓扑空间,
则 $E=\Top(X)$ 的子拓扑斯 $E'$ 是一个非“空”而没有点的拓扑斯.

\item 设 $X'$ 是 $X$ 的一个子集,
$E_{X'}$ 是经 \hyperref[IV.exer9.1.8]{9.1.8} a)
的构造与 $X'$ 对应的 $E$ 的子拓扑斯.
证明若 $X'$ 稠密, 则 $E_{X'}$ 包含 $E'$.
证明若 $X$ 清醒, 且 $X'$ 和 $X''$ 是 $X$ 的两个清醒子集,
则交 $X'\cap X''$ 清醒,
并有 $E_{X'\cap X''}\subset E_{X'}\cap E_{X''}$,
而且这个包含可以是严格的
\footnote[1]{不过, 当 $X$ 局部 Noether 时, 见下文
\hyperref[IV.exer9.1.11]{9.1.11} b).}.
(取一个非空 Hausdorff 空间 $X$, 使它有两个不交的稠密子集 $X'$, $X''$.)
证明上述包含为等式,
当且仅当子拓扑斯 $E_{X'}\cap E_{X''}$ 有足够多的点.
证明若 $X'$ 或 $X''$ 是 $X$ 的\emph{局部闭}子集, 则情况正是如此.

\item 设 $X'$ 是 $X$ 的一个子集, $X''$ 是 $X'$ 的补集.
$E$ 的子拓扑斯 $F$ 包含对 $x\in X''$ 的全部 $E_{\{x\}}$,
当且仅当它包含 $E_{X''}$
(\cf \hyperref[IV.exer9.1.8]{9.1.8} b)).
由此得出: 若 $X''$ 是 $X$ 的局部闭子集的并
(例如 $X''$ 的点均为闭点),
则 $E$ 的所有与 $E_{X'}$ 的交为空子拓扑斯 $\Top(∅)$ 的子拓扑斯 $F$ 所成集合,
\emph{只有当} $E_{X'}\cap E_{X''}$ 为空子拓扑斯时,
才有最大元素
(它此时才当得起 $E_{X'}$ 在 $E$ 中的\emph{弱补}之名;
\cf \hyperref[IV.exer9.1.13]{9.1.13})
(若 $X\neq∅$, $X'$ 和 $X''$ 在 $X$ 中都稠密,
且 $X$, $X'$, $X''$ 均清醒, 则这个条件不满足; \cf c)).
证明当此条件不满足时,
在 $E$ 的子拓扑斯所成的有序集中,
两个元素的下确界对任意 $\Sup$ 不满足分配律.
(编者不知道它是否对有限 $\Sup$ 满足分配律;
不过 \cf \hyperref[IV.exer9.1.11]{9.1.11} e)
和 \hyperref[IV.exer9.1.12]{9.1.12} d).)
\end{enumerate}
\end{enonce*}

\begin{enonce*}[remark]{习题 9.1.11}\etiquette[9.1.11]{IV.exer9.1.11}
(\emph{局部 Noether 拓扑斯的子拓扑斯. Noether 空间的情形}).
\begin{enumerate}
\renewcommand{\theenumi}{\alph{enumi}}
\renewcommand{\labelenumi}{\theenumi)}
\item 设 $f:E'\rightarrow E$ 是一个拓扑斯嵌入.
证明若 $X$ 是 $E$ 的一个预 Noether 对象
(\ie (\sisi{VI 1.}{\hyperrefext[VI][VI.def1.30]{1.30}})
$X$ 的每个递增子对象序列\pageoriginale 都稳定),
则 $X'=f^{\ast}(X)$ 是 $E'$ 的一个预 Noether 对象.
(引入诱导拓扑斯 ${E'}_{/X'}$ 和 $E_{/X}$,
并利用 \hyperref[IV.exer9.1.6]{9.1.6} f),
先归结到 $X$ 是 $E$ 的终对象 $e$ 的情形,
于是 $X'$ 是 $E'$ 的终对象 $e'$.
再注意, $e'$ 的每个递增子对象序列,
应用 $f_{\ast}$ 后都定义 $f_{\ast}(e')=e$ 的一个递增子对象序列.)
由此得出: 若 $E$ 有一个由预 Noether 对象组成的生成族
(\resp 若 $E$ 是 Noether 拓扑斯,
\resp 局部 Noether 拓扑斯 (\sisi{VI 2.)}{\hyperrefext[VI][VI.def2.11]{2.11}})),
则 $E'$ 满足同样的条件.
(利用 \hyperref[IV.exer9.1.9]{9.1.9} b).)

\item 设 $E$ 是一个局部 Noether 拓扑斯
(\sisi{VI 2.}{\hyperrefext[VI][VI.def2.11]{2.11}}).
证明 $E$ 的每个子拓扑斯都形如 $E_Z$
(\hyperref[IV.exer9.1.8]{9.1.8} a)),
其中 $Z$ 是 $\Ppoint(E)$ 的某个适当全子范畴.
(利用 a) 和 \textsc{Deligne} 关于每个局部 Noether 拓扑斯都有足够多点的定理.)
由此得出: 若 $E$ 形如 $\Top(X)$,
其中 $X$ 是一个清醒的局部 Noether 拓扑空间,
则 $Z\mapsto E_Z$ 是从 $X$ 的清醒子空间所成的有序集,
到 $E$ 的子拓扑斯所成的有序集的一个同构.
(利用 \hyperref[IV.exer9.1.8]{9.1.8} e).)
清醒子空间 $Z_i$ 的交必清醒,
并定义子拓扑斯 $E_{Z_i}$ 的交
(这与一般情形 (\hyperref[IV.exer9.1.10]{9.1.10} c) 可能发生的情况不同).

\item 设 $X$ 是一个清醒的局部 Noether 空间,
$X'$ 是 $X$ 的一个子集, $X''$ 是其补集.
证明以下条件等价:
\begin{enumerate}
\renewcommand{\theenumii}{\roman{enumii}}
\renewcommand{\labelenumii}{(\theenumii)}
\item $X'$ 可构造.

\item $X'$ 和 $X''$ 清醒.

\item $X'$ 清醒, 且 $E=\Top(X)$ 的子拓扑斯
$E_{X'}\approxeq\Top(X')$ 有一个“互补”子拓扑斯
(\cf \sisi{\hyperref[IV.exer9.1.13]{9.1.13} e)}{\hyperref[IV.exer9.1.13]{9.1.13}}).

\item $E=\Top(X)$ 的子拓扑斯 $E_{X'}\approx\Top(X')$
和 $E_{X''}\approx\Top(X'')$ 的交是 $E$ 的“空”子拓扑斯.
\end{enumerate}

证明: 若 $E'$ 和 $E''$ 是 $E$ 的两个子拓扑斯,
则它们互为弱补, 当且仅当它们互补
(\hyperref[IV.exer9.1.13]{9.1.13}).

\item 设\pageoriginale $X$ 是一个清醒 Noether 空间.
证明以下条件等价:
\begin{itemize}
\item[(i)] $X$ 的每个清醒子空间 $X'$ 都可构造,
\ie 根据 c), 它的补集清醒.

\item[(i bis)] 对 $X$ 的每个极大点 $x$,
$X-\{x\}$ 清醒, \ie $\{x\}$ 可构造.

\item[(ii)] $X$ 有限.

\item[(iii)] $\Top(X)$ 的每个子拓扑斯都有一个互补子拓扑斯
(\cf \hyperref[IV.exer9.1.13]{9.1.13}).

\item[(iv)] 在 $\Top(X)$ 的子拓扑斯所成的有序集中,
两个元素的下确界对任意 $\Sup$ 满足分配律.
\end{itemize}
(利用 \hyperref[IV.exer9.1.10]{9.1.10} d) 的论证,
并结合 b) 中证明的交的公式.)

\item 设 $X$ 是一个局部 Noether 空间.
证明在 $\Top(X)$ 的子拓扑斯所成的有序集中,
两个元素的下确界对\emph{有限} $\Sup$ 满足分配律.
(归结到 $X$ 清醒的情形, 并利用 b).)

\item 对清醒的局部 Noether 空间 $X$ 的每个清醒子空间 $Z$,
设 $T_Z$ 是 $X$ 的开子集范畴 $\Ouv(X)$ 上的拓扑,
其中包含族 $U_i\rightarrow U$ 覆盖,
当且仅当 $U\cap Z=\bigcup_i U_i\cap Z$.
证明 $Z\mapsto T_Z$ 在 $X$ 的清醒子空间 $Z$ 所成的集合,
与 $\Ouv(X)$ 上比典范拓扑 $T_X$ 更细的拓扑所成的集合之间,
建立一个 (反转包含关系的) 双射.
(利用 b) 和 \hyperref[IV.rem9.1.2]{9.1.2} c).)

\item 设 $X$ 是一个清醒拓扑空间.
$X$ 的一个子集 $X'$ 清醒, 当且仅当其补集为局部闭子集的并;
若 $X$ 局部 Noether, 这也就是说 $X'$ 是“投射可构造的” (EGA IV 1.9.4).
证明 $X$ 上存在一个拓扑 $T_{\cons'}$,
使 $X_{\cons'}=(X,T_{\cons'})$ 的闭子集是 $X$ 的清醒子集.
此后假设 $X$ 局部 Noether,
于是 $X_{\cons'}$ 正是 EGA IV 9.1.13 中的空间 $X_{\cons}$\pageoriginale.
由 b) 和 c) 得出: $\Top(X)$ 的子拓扑斯所成的有序集 $\Sigma$,
同构于 $X_{\cons}$ 的闭子集所成的有序集;
在这一对应下, 有补的子拓扑斯对应于 $X_{\cons}$ 中既开又闭的子集.
证明 \sisi{$X$}{$X_{cons}$} 局部紧,
且它紧, 当且仅当 $X$ 是 Noether 空间,
\ie 拓扑斯 \sisi{$\Top(E)$}{$\Top(X)$} 是 Noether 的.

\item 设 $E$ 是一个局部 Noether 拓扑斯.
假设它具有 \hyperref[IV.subsubsec9.1.14]{9.1.14} b) 中考虑的性质
(这个条件或许总是成立\ldots).
证明典范函子 $\Ppoint(E')\rightarrow\Ppoint(E)$ 的本质像,
其中 $E'$ 遍历 $E$ 的子拓扑斯集合 $\Sigma$,
定义了 $E$ 的点的同构类所成的集合 $X=\Point(E)$ 的一族子集,
这族子集正是 $X$ \emph{上}某个\emph{拓扑} $T_{\cons}$ 的闭子集族;
由此得到 $E$ 的子拓扑斯集合 $\Sigma$
与 $X_{\cons}$ 的闭子集集合之间的一个有序集同构.
由此得出: 在 $\Sigma$ 中, 两个元素的 $\Sup$ 对\emph{任意}下确界满足分配律.
证明不一定为 $𝒰$-小的 $X_{\cons}$
有一个 $𝒰$-小的相伴清醒空间 (\hyperref[IV.subsubsec4.2.1]{4.2.1}).
\end{enumerate}
\end{enonce*}

\begin{enonce*}[remark]{习题 9.1.12}\etiquette[9.1.12]{IV.exer9.1.12}
(\emph{有限拓扑斯}).
若拓扑斯 $E$ 等价于形如 $\widehat{C}$ 的拓扑斯,
其中 $C$ 是有限范畴, 就称它\emph{有限}.
\begin{enumerate}
\renewcommand{\theenumi}{\alph{enumi}}
\renewcommand{\labelenumi}{\theenumi)}
\item (对应字典). 设 $𝒱$ 是一个满足 $𝒰\in𝒱$ 的宇宙.
设 (Karfiness) 是范畴 $(\Cat)_{𝒱}$ 的如下全 2-子范畴定义的 2-范畴:
其对象是属于 $𝒱$、\emph{为 Karoubi 且等价于有限范畴}的范畴;
设 (Topfin) 是 $(𝒱\text{-}𝒰\text{-}\Top)$
(\hyperref[IV.subsubsec3.3.1]{3.3.1}) 的如下全 2-子范畴定义的 2-范畴:
其对象是 $\in𝒱$ 的\emph{有限} $𝒰$-拓扑斯.
证明存在互为拟逆的 2-等价:
\hypertarget{sga4.exp.IV.equation.9.1.12.1}{}\begin{align*}
C & \longmapsto \widehat{C}: (\text{Karfiness})\longrightarrow
(\text{Topfin})\tag{9.1.12.1}\\
E & \longmapsto \Ppoint(E): (\text{Topfin})\longrightarrow
(\text{Karfiness}).\tag{9.1.12.2.}
\end{align*}\etiquette[9.1.12.2]{IV.eq9.1.12.2}
(利用 \hyperref[IV.exer7.6]{7.6} h), e).)

\item 证明\pageoriginale 有限拓扑斯是 Noether 的
\sisi{(VI 2)}{\hyperrefext[VI][VI.def2.11]{2.11}}.

\item 证明有限拓扑斯 $E$ 的每个子拓扑斯都是有限拓扑斯.
更确切地, 若 $E$ 等价于 $\widehat{C}$,
其中 $C$ 是有限 Karoubi 范畴
(或更一般地, 等价于某个有限范畴),
证明通过将每个 $C'$ 对应到 $f_{\ast}:\widehat{C}'\rightarrow \widehat{C}$ 的本质像,
其中 $f:C'\rightarrow C$ 是包含函子,
便在 $E$ 的子拓扑斯 $E'$ 所成的集合,
与 $C$ 中为 Karoubi 的
(或等价地, 对 $C$ 中的直因子稳定的) 严格全子范畴 $C'$ 所成的集合之间,
得到一个序同构.
(利用 \hyperref[IV.subsec5.6]{5.6} 验证 $f_{\ast}$ 全忠实,
并利用 b) 及 \hyperref[IV.exer9.1.11]{9.1.11} b)
证明 $E$ 的每个子拓扑斯 $E'$ 均可按所述方式得到.)

\item 设 $E$ 是一个有限拓扑斯.
在 $E$ 的点的同构类所成的有限集合 $X=\Point(E)$ 上构造一个拓扑,
使它成为清醒拓扑空间,
且 $E$ 的子拓扑斯所成的有序集典范同构于 $X$ 的闭子集集合.
(取 $X$ 上的“可构造拓扑” $\Cc_{\cons}$,
它由序关系 $(x\leq y)\Leftrightarrow$
($x$ 是 $y$ 的一个直因子) 定义,
对于这个拓扑, 点 $x$ 的闭包由满足 $y\leq x$ 的 $y$ 组成.)
特别地, $E$ 的子拓扑斯集合 $\Sigma$ 有限,
且二元下确界运算 $\wedge$ 对任意 $\Sup$ 满足分配律.

\item 设 $C$ 是一个等价于有限范畴的范畴.
证明对 $C$ 的每个全子范畴 $C'$,
存在 $C$ 上的一个拓扑 $T_{C'}$,
使族 $(X_i\rightarrow X)_{i\in I}$ 覆盖,
当且仅当对每个 $Y\in\Ob C'$,
映射族 $\Hom(Y,X_i)\rightarrow\Hom(Y,X)$ $(i\in I)$ 满射.
证明 $(C,T_{C'})$ 上的层范畴等价于 $\widehat{C}'$,
且 $C'\mapsto T_{C'}$ 是如下两个集合之间的一个 (反转序关系的) 双射:
$C$ 中为 Karoubi 的
(或等价地, 对 $C$ 中的直因子稳定的) 严格全子范畴 $C'$ 所成的集合,
以及范畴 $C$ 上的拓扑 (\hyperrefext[II][II.def1.1]{1.1}) 所成的集合.
(利用 $c$) 和 \hyperref[IV.rem9.1.2]{9.1.2} c).)\pageoriginale
特别地, 若景 $C$ 的底层范畴等价于有限范畴,
则拓扑斯 $C^{\sim}$ 是有限拓扑斯.

\item 设 $f:C'\rightarrow C$ 是一个 $𝒰$-范畴函子,
其中 $C$ 等价于有限范畴.
证明拓扑斯态射 $\widehat{f}:\widehat{C}'\rightarrow \widehat{C}$
(\hyperref[IV.subsec4.6]{4.6}) 保守,
当且仅当 $C$ 的每个对象都同构于 $f$ 的像中某个对象的直因子.
(归结到 $f$ 是一个包含 $f:C'\hookrightarrow C$ 的情形,
其中 $C'$ 如 c) 所述.)

\item 设 $C$ 是一个有两个对象 $e$ (终对象) 和 $a$ 的范畴,
除恒等箭头外, 还有三个箭头 $f:a\rightarrow e$,
$g:e\rightarrow a$, $p=gf$ (因此 $p^2=p$).
证明 $E=\widehat{C}$ 除子拓扑斯 $E$ 和 $\Top(∅)$ 之外,
恰好只有一个子拓扑斯 $E'$, 即 $E'=\widehat{\{e\}}$;
因此, 对 $E$ 的每个满足 $E''\neq\Top(∅)$ 的子拓扑斯 $E''$,
有 $E'\cap E''\neq\Top(∅)$.
\end{enumerate}
\end{enonce*}

% END INCLUDED SGA4_English_current_buildable_source_20260906/source/04/drafts/SGA4_IV_LINES_5779_6107_English.texfrag


% BEGIN INCLUDED SGA4_English_current_buildable_source_20260906/source/04/drafts/SGA4_IV_LINES_6108_6401_English.texfrag
\begin{enonce*}[remark]{习题 9.1.13}\etiquette[9.1.13]{IV.exer9.1.13}
(\emph{拓扑斯的互补子拓扑斯}).

设 $E$ 是一个拓扑斯, $E'$ 和 $E''$ 是两个子拓扑斯.
若 $E'\cap E''=\Top(∅)$ 且 $E'\vee E''=E$
(其中 $\vee$ 表示 $E$ 的子拓扑斯所成有序集中的 $\Sup$),
则称 $E'$ 和 $E''$ \emph{互补}.

假设在 $E$ 的子拓扑斯集合 $\Sigma$ 中,
两个元素的下确界对有限上确界满足分配律.

\begin{enumerate}
\renewcommand{\theenumi}{\alph{enumi}}
\renewcommand{\labelenumi}{\theenumi)}
\item 若 $E'$ 和 $E''$ 互补,
则 $E''$ 是满足 $E'\cap E''_1=\Top(∅)$ 的子拓扑斯 $E''_1$ 中最大的一个
(\ie 用 \hyperref[IV.exer9.1.10]{9.1.10} d) 的术语来说,
$E''$ 是 $E'$ 的一个弱补).

\item 证明 $E$ 的以下条件等价:
\begin{itemize}
\item[(i)] 若 $E'$ 和 $E''$ 是 $E$ 的两个子拓扑斯,
且 $E''$ 是 $E'$ 的一个弱补,
则 $E''$ 是 $E'$ 的一个补.

\item[(ii)] 对 $E$ 的每个满足 $E'\neq E$ 的子拓扑斯 $E'$,
存在 $E$ 的一个子拓扑斯 $E''$,
使 $E''\neq\Top(∅)$ 且 $E'\cap E''=\Top(∅)$.

\item[(i bis)] 若 $E'$ 和 $E''$ 是 $E$ 的两个子拓扑斯,
且 $E''$ 在满足 $E'\cap E''_1=\Top(∅)$ 的子拓扑斯 $E''_1$ 中极大,
则 $E''$ 是 $E'$ 的一个补.
\end{itemize}

证明即使 $E$ 是有限拓扑斯 (\hyperref[IV.exer9.1.12]{9.1.12}),
这些\pageoriginale 条件也未必成立
(\hyperref[IV.exer9.1.12]{9.1.12} g)).
不过, 若 $E$ 形如 $\Top(X)$,
其中 $X$ 是一个局部 Noether 空间, 则它们成立.
(利用 \hyperref[IV.exer9.1.11]{9.1.11} b).)

\item 若 $E$ 的子拓扑斯 $E'$ 有一个补,
就称它\emph{有补} (比较 \hyperref[IV.exer9.1.12]{9.1.12} c)).
证明有补子拓扑斯所成的集合对有限下确界、有限上确界及取补稳定,
并且取补将下确界变为上确界, 将上确界变为下确界.

\item 陈述 a), b), c) 中各项观察的对偶
(它们实际上是关于抽象有序集 $\Sigma$ 的显然的一般陈述;
注意, 关于 $\Sigma$ 的假设本身自对偶).
\end{enumerate}
\end{enonce*}

\setcounter{subsubsection}{13}
\subsubsection{}\etiquette[9.1.14]{IV.subsubsec9.1.14}
\emph{开放问题}\nde{问题 a) 和 d) 的肯定回答见 [17].}.
设 $E$ 是一个拓扑斯.

\begin{itemize}
\item[a)] 在 $E$ 的子拓扑斯所成的集合中,
两个元素的下确界对有限上确界是否满足分配律?

\item[b)] 设 $E'$ 和 $E''$ 是 $E$ 的两个子拓扑斯,
满足 $E=E'\vee E''$.
那么 $E$ 的每个点是否都同构于 $E'$ 或 $E''$ 的某个点的像?

注意, 若对 $E$ 及其所有子拓扑斯, b) 的回答都是肯定的,
则至少对于有足够多点的子拓扑斯, a) 的回答是肯定的
(因而若 $E$ 局部 Noether, 则对所有子拓扑斯都如此
(\hyperref[IV.exer9.1.11]{9.1.11} b)).
不过, 对所有 Noether 拓扑斯, a) 的回答
(更不用说 b) 的回答) 仍未知.
但是, 若 $E$ 形如 $\Top(X)$,
其中 $X$ 是局部 Noether 空间 (\hyperref[IV.exer9.1.11]{9.1.11}),
或 $E$ 是有限拓扑斯 (\hyperref[IV.exer9.1.12]{9.1.12}),
则 a) 和 b) 的回答都是肯定的.

注意, a) 可以重述如下 (应用于 $E$ 的所有子拓扑斯 $F$):

\item[$\rm a'$)] 若 $E'$ 和 $E''$ 是拓扑斯 $F$ 的两个子拓扑斯,
使典范态射 (\cf \hyperref[IV.exer8.7]{8.7} b))
\begin{equation*}
E'\amalg E''\longrightarrow F\tag{9.1.14.1}
\end{equation*}\etiquette[9.1.14.1]{IV.eq9.1.14.1}
保守 (这也就是说 $F=E'\vee E''$
(\hyperref[IV.exer9.1.7.2]{9.1.7.2} e)),
则经 $F$ 的子拓扑斯 $F_1$ 的 2-基变换 $F_1\hookrightarrow F$\pageoriginale 后仍然保守.

b) 也有类似的重述 $\rm b'$),
只须用一个\emph{点}拓扑斯 $F_1$ 作基变换 $F_1\rightarrow F$
(\hyperref[IV.subsec2.2]{2.2}).
因此, 若以下问题的回答是肯定的
(其中 $F$ 是 $E$ 的任意子拓扑斯),
则 a) 和 b) 的回答都是肯定的:

\item[c)] 若 $E'$ 和 $E''$ 是拓扑斯 $F$ 的两个子拓扑斯,
其上确界为 $F$,
\ie 拓扑斯态射 \eqhyperref[IV.eq9.1.14.1]{9.1.14.1} 保守,
则这个态射是否甚至\emph{普遍保守},
\ie 在每次 2-基变换 $F_1\rightarrow F$ 后都保持保守?

由 \textsc{Deligne} 关于拓扑斯的 2-纤维积存在性的结果
(见 [16], 命题 4.47), 这个陈述是有意义的.
不过, 见 \hyperref[IV.exer9.1.7.2]{9.1.7.2} d) 中的例子.

另一方面, \hyperref[IV.exer9.1.11]{9.1.11} h) 表明,
在局部 Noether 情形中, b) 的肯定回答将给出以下问题的肯定回答:

\item[d)] 在 $E$ 的子拓扑斯所成的有序集 $\Sigma$ 中,
两个元素的上确界对\emph{任意} (不一定有限) 下确界是否满足分配律?

这也就是说, $\Sigma$ 可以解释为某个适当拓扑斯 $E_{\cons'}$
的“闭子拓扑斯”所成的有序集
(即 $\Sigma^{0^{\sim}}$,
其中将 $\Sigma$ 的反向有序集 $\Sigma^0$ 视为范畴,
并赋予其典范拓扑, \cf \hyperref[IV.exer7.8]{7.8} b));
这个拓扑斯由终对象的子对象生成, 因而与通常拓扑空间非常接近.
还有待研究这个拓扑斯 $E_{\cons'}$,
可从例 \hyperref[IV.exer9.1.11]{9.1.11} g)
和 \hyperref[IV.exer9.1.12]{9.1.12} d) 获得启发;
特别要确定当 $E$ 为 Noether 时它是否也为 Noether.
这等于对 $E$ 的子拓扑斯 $F$ 提出以下问题,
它独立于 d) 也有意义:

\item[e)] 设 $F$ 是一个 Noether 拓扑斯,
$(E_i)_{i\in I}$ 是 $F$ 的一族子拓扑斯, 其交为 $\Top(∅)$.
是否存在一个具有相同性质的有限子族?

在一个相当不同的方向上,
还应回顾 \hyperref[IV.exer9.1.6]{9.1.6} e) 中提出的
问题\pageoriginale, 它值得澄清.
\end{itemize}

\subsection{开子拓扑斯的情形}\etiquette[9.2]{IV.subsec9.2}

\subsubsection{}\etiquette[9.2.1]{IV.subsubsec9.2.1}

设 $U$ 是拓扑斯 $E$ 的一个开子对象,
\ie $E$ 的终对象 $e$ 的一个子对象 (\hyperref[IV.def8.3]{8.3}).
考虑局部化态射 (\hyperref[IV.subsec5.2]{5.2})
\begin{equation*}
j:E_{/U}\longrightarrow E{\quoi}.\tag{9.2.2}
\end{equation*}\etiquette[9.2.2]{IV.eq9.2.2}
由于 $U\rightarrow e$ 是单态射,
函子 $j_{!}$ (可解释为遗忘函子) 全忠实,
从而它的双伴随 $j_{\ast}$ 也全忠实
(\sisi{1.5.7 a)}{\hyperrefext[I][I.rem5.7]{5.7}}).
换言之, \emph{\hyperref[IV.eq9.2.2]{(9.2.2)} 中
与拓扑斯 $E$ 的一个开子对象 $U$ 相应的局部化态射 $j$ 是一个拓扑斯嵌入}
(\hyperref[IV.def9.1.1]{9.1.1} a)).

若拓扑斯嵌入 $F\rightarrow E$ 同构于由 $E$ 的某个开子对象 $U$
定义的拓扑斯嵌入, 就称它是\emph{开嵌入}.
若 $E$ 的子拓扑斯 $E'$ (\hyperref[IV.def9.1.1]{9.1.1} b))
由 $E$ 的某个开子对象定义,
\ie 典范包含态射 $E'\rightarrow E$ 是开嵌入,
就称它是\emph{开子拓扑斯}.
注意, 与开嵌入 $j:E'\rightarrow E$ 相应的 $E$ 的开子对象 $U$
唯一确定为 $e$ 的子对象 $j_{!}(e')$,
其中 $e'$ 表示 $E'$ 的终对象.
这样就得到 $E$ 的开子对象 $U$ 与 $E$ 的开子拓扑斯之间的一一对应.

\begin{enonce*}[remark]{注 9.2.3}\etiquette[9.2.3]{IV.rem9.2.3}
按照上述定义, 与 $E$ 的开子对象 $U$ 相应的 $E$ 的开子拓扑斯,
是 $E$ 的如下严格全子范畴:
它是正像函子
$$
j_{\ast}:E_{/U}\longrightarrow E.
$$
的本质像.
不应将 $E$ 的这个子范畴与 $j_{!}$ 的本质像所成的严格全子范畴混淆;
后者由 $E$ 中使结构态射 $X\rightarrow e$ 经过 $U$ 分解的对象 $X$ 组成.
这两个子范畴显然互相确定, 并且典范等价.
因此, 某些作者可能倾向于把
(而且确实把) $j_{!}$ 的本质像所成的严格全子范畴
称为“由 $U$ 定义的开子拓扑斯”,
因为用 $U$ 描述它比描述 $j_{\ast}$ 的本质像\pageoriginale 更简单.
只要不打算研究开子拓扑斯之外的其他种类子拓扑斯,
这个术语就不会造成困难.
我们的情形并非如此, 因而不采用上述作者的用法.

为方便读者, 我们概括开嵌入最重要的特殊性质:
\end{enonce*}

\begin{enonce*}{命题 9.2.4}\etiquette[9.2.4]{IV.prop9.2.4}
设
$$
j:\mathcal{U}\longrightarrow E
$$
是一个开拓扑斯嵌入 (\hyperref[IV.subsubsec9.2.1]{9.2.1}).
则 $j^{\ast}$ 的左伴随函子 $j_{!}$ 存在,
所以有三个伴随函子的序列
\[
\xymatrix@R=.35cm@C=.5cm{
           & j_!,\quad j^{\ast}, \quad j_{\ast} : &\\
&&\\
\mathcal{U} \ar @<-16pt>[rr]^{j_{\ast}} \ar
           @<14pt>[rr]^{j_{!}}&            &\ar @<1pt>[ll]_{j^{\ast}} E
}
\]

此外:
\begin{enumerate}
\renewcommand{\theenumi}{\alph{enumi}}
\renewcommand{\labelenumi}{\theenumi)}
\item 函子 $j_{!}$ 和 $j_{\ast}$ 全忠实.

\item 函子 $j_{!}$ 与纤维积、以小集合 $I\neq\phi$ 为指标的积,
以及以小的余滤过范畴为指标的投射极限交换
(\hyperrefext[I][I.def2.7]{2.7}).

\item 对 $E$ 的每个对象 $X$, 伴随态射
$$
j_{!}j^{\ast}(X)\longrightarrow X
$$
是单态射.
\end{enumerate}
\end{enonce*}

\noindent \emph{证明}.
可假设 $j$ 是由 $E$ 的一个开子对象 $U$ 定义的局部化态射.
此时 a) 只是为便于引用而重述;
b) 来自将 $j_{!}$ 解释为遗忘函子,
并利用 $U\rightarrow e$ 是单态射这一事实.
最后, 注意到所考虑的态射是包含 $U\rightarrow e$
沿 $X\rightarrow e$ 的基变换,
而单态射经基变换仍为单态射, c) 便显然成立.

\begin{enonce*}[remark]{注 9.2.4.1}\etiquette[9.2.4.1]{IV.rem9.2.4.1}
存在一些拓扑斯嵌入 $j:E'\rightarrow E$,
其 $j_{!}$ 存在, 却不与两个对象的积或余滤过投射极限交换
(习题 \hyperref[IV.exer9.5.9]{9.5.9} c)).
\end{enonce*}

\setcounter{subsubsection}{4}
\subsubsection{}\etiquette[9.2.5]{IV.subsubsec9.2.5}
现在\pageoriginale 设
$$
g:E'\longrightarrow E
$$
是一个任意拓扑斯态射,
$U'=g^{\ast}(U)$ 是 $E$ 的开子对象 $U$ 的逆像.
我们在 \hyperref[IV.prop5.11]{5.11} 已看到, 拓扑斯图
\[
\xymatrix@C=2cm@R=1.5cm{
E'/_{U'} \ar[d]_{g_U} \ar[r]^{j'}& E' \ar[d]^{g}\\
E/_{U}  \ar[r]^{j} & E }
\]
是 2-Cartesian 的.
特别可知, $E$ 的开子拓扑斯在 $g$ 下的逆像
(\cf \hyperref[IV.exer9.1.6]{9.1.6} d)) 是 $E'$ 的开子拓扑斯;
或等价地, 在 $𝒰$-拓扑斯
(属于给定宇宙 $𝒱$) 的 2-范畴中,
开拓扑斯嵌入 (\hyperref[IV.subsubsec9.2.1]{9.2.1}) 的概念对 2-基变换稳定.

还应注意, 拓扑斯态射 $g:E'\rightarrow E$
在相差同构的意义下经过 $j:E_{/U}\rightarrow E$ 分解,
当且仅当 $j':E'_{/U'}\rightarrow E'$ 是范畴等价,
这也就是说, 包含 $U'\hookrightarrow e'$ 是同构
(其中 $e'$ 为 $E'$ 的终对象).
这通过简单刻画 \loccit 中考虑的函子的本质像,
使 \hyperref[IV.prop9.1.4]{9.1.4} 更为明确.

\subsubsection{}\etiquette[9.2.6]{IV.subsubsec9.2.6}
将 \hyperref[IV.subsubsec9.2.5]{9.2.5}
应用于 $E'$ 形如 $E_{/V}$ 的情形,
其中 $V$ 是 $E$ 的第二个开子对象, $g$ 是局部化态射.
此时 $U'=U\cap V$,
$E_{/U}$ 和 $E_{/V}$ 在 $E$ 上的 2-Cartesian 积可与 $E_{/U\cap V}$ 等同.
由此可知, \emph{$E$ 的开子拓扑斯的有限交是 $E$ 的开子拓扑斯};
更确切地, 映射
\begin{equation*}
U\longmapsto \text{ 由 $U$ 定义的 $E$ 的子拓扑斯}
\tag{9.2.6.1}
\end{equation*}\etiquette[9.2.6.1]{IV.eq9.2.6.1}
与有限交交换.

\begin{enonce*}[remark]{习题 9.2.7}\etiquette[9.2.7]{IV.exer9.2.7}
证明映射 \eqhyperref[IV.eq9.2.6.1]{9.2.6.1}
也与\emph{任意} $\Sup$ 交换.
(利用 \hyperref[IV.exer9.1.7.2]{9.1.7.2} e).)
\end{enonce*}

% END INCLUDED SGA4_English_current_buildable_source_20260906/source/04/drafts/SGA4_IV_LINES_6108_6401_English.texfrag


% BEGIN INCLUDED SGA4_English_current_buildable_source_20260906/source/04/drafts/SGA4_IV_LINES_6402_6597_English.texfrag
\subsection{与开子拓扑斯互补的闭子拓扑斯的构造}\etiquette[9.3]{IV.subsec9.3}\pageoriginale

\subsubsection{}\etiquette[9.3.1]{IV.subsubsec9.3.1}
设 $\mathsf{T}$ 是一个拓扑空间, $U$ 是 $\mathsf{T}$ 的一个开子集,
于是 $\Top(U)$ 可与 $\Top(\mathsf{T})$ 的一个开子拓扑斯等同
(沿用 \hyperref[IV.subsec2.1]{2.1} 的记号).
设 $Y$ 是 $\mathsf{T}$ 中与 $U$ 互补的闭拓扑子空间.
在相差等价的意义下, $\Top(Y)$ 可视为 $\Top(\mathsf{T})$ 的一个子拓扑斯,
即由 $\Top(\mathsf{T})$ 中限制在 $U$ 上为
\sisi{$U$ 的终拓扑斯}{$U$ 的终对象}的对象 $F$ 组成的子拓扑斯.
只要有一个拓扑斯及其开子对象 $U$,
关于 $E=\Top(\mathsf{T})$ 的子范畴的这个描述就有意义,
而且我们将看到它总给出 $E$ 的一个子拓扑斯.
此外, 在最初考虑的特殊情形中,
直接可见 (采用习题 \hyperref[IV.exer9.1.13]{9.1.13} 引入的术语),
$\Top(Y)$ 和 $\Top(U)$ 是互补子拓扑斯
(利用 \hyperref[IV.subsubsec9.2.5]{9.2.5}
和 \hyperref[IV.exer9.1.7.2]{9.1.7.2} e)).
一般情形也如此, 而且我们将看到这个性质唯一刻画所讨论的子拓扑斯.
因此, 无论从何种意义说,
它都当得起与给定开子对象 $U$ (或由 $U$ 定义的开子拓扑斯) 互补的闭子拓扑斯之名.
本节给出此拓扑斯 $\mathcal{F}$ 和逆像函子
$i^{\ast}:E\rightarrow \mathcal{F}$ 的详细构造,
而第 \hyperref[IV.subsec9.4]{9.4} 节讨论它的初步性质.

\subsubsection{}\etiquette[9.3.2]{IV.subsubsec9.3.2}
因此, 设 $E$ 是一个拓扑斯, $U$ 是 $E$ 的一个开子对象,
如 \hyperref[IV.subsubsec9.2.1]{9.2.1} 所述,
并考虑局部化态射
\begin{equation*}
j:\mathcal{U}=E_{/U}\longrightarrow E{\quoi},\tag{9.3.2.1}
\end{equation*}\etiquette[9.3.2.1]{IV.eq9.3.2.1}
它是一个开嵌入.

对 $E$ 的每个对象 $X$, 置
\begin{equation*}
X_{\complement U}=U\coprod_{U\times X}
X{\quoi},\tag{9.3.2.2}
\end{equation*}\etiquette[9.3.2.2]{IV.eq9.3.2.2}
其中纤维余积关于典范投射
$$
\pr_1:U\times X\longrightarrow U,\quad \pr_2:U\times
X\longrightarrow X.
$$
形成. 因而 $E$ 中有一个关于 $X$ 具有函子性的
\emph{余 Cartesian} 图 (\hyperrefext[I][I.sec10]{10}):\pageoriginale
\[\vcenter{
\xymatrix@C=1.75cm@R=1.5cm{
U\times X \ar[r]^{\pr_2} \ar[d]_{\pr_1}& X \ar[d]\\
U \ar[r]        & X_{\complement U} }}\leqno{(9.3.2.3)}
\]\etiquette[9.3.2.3]{IV.eq9.3.2.3}
注意, 在 \hyperref[IV.subsubsec9.3.1]{9.3.1} 考虑的例子中,
$X_{\complement U}$ 典范同构于 $i_{\ast}i^{\ast}(X)$
(其中 $i:Y\rightarrow T$ 为包含),
且 $X\rightarrow X_{\complement U}$ 可与伴随态射等同.

\begin{enonce*}{命题 9.3.3}\etiquette[9.3.3]{IV.prop9.3.3}
沿用上述记号, 有以下陈述:
\begin{enumerate}
\renewcommand{\theenumi}{\alph{enumi}}
\renewcommand{\labelenumi}{\theenumi)}
\item 对 $E$ 的每个对象 $X$, 以下条件等价:
\begin{itemize}
\item[(i)] 存在 $E$ 的一个对象 $Y$ 以及一个同构 $X\simeq Y_{\complement U}$.

\item[(ii)] 典范态射 $X\xrightarrow{m}X_{\complement U}$ 是同构.

\item[(iii)] 典范态射 $\pr_2:X\times U\rightarrow U$ 是同构
(\ie $j^{\ast}(X)$ 是 $E_{/U}$ 的一个终对象).
\end{itemize}

\item 对 $E$ 中的每个态射
$$
m:X\longrightarrow X'
$$
以下条件等价:
\begin{itemize}
\item[${\rm (i')}$] 图
\[
\xymatrix@C=2cm@R=1.5cm{
X \ar[r]^{m} \ar[d]& X'\ar[d]\\
X_{\complement U} \ar[r]& X'_{\complement U} }
\]
是 Cartesian 的.

\item[${\rm (ii')}$] 态射 $m\times\id_U:X\times U\rightarrow X'\times U$
是同构, \ie $j^{\ast}(m)$ 是同构.
\end{itemize}

\item 从 $E$ 到 $E$ 的函子
$$
X\longrightarrow X_{\complement U}
$$
左正合, 将满态射族映为满态射族,
并且与滤过归纳极限及\pageoriginale 纤维余积交换.
\end{enumerate}
\end{enonce*}

\noindent \emph{证明}.
先在 $E$ 形如 $\widehat{C}$ 的情形中证明此命题,
其中 $C$ 是一个小范畴.
由 \hyperref[I.prop3.1]{I.3.1},
这归结到 $E$ 是“点”拓扑斯 $(\Ens)$ 的情形, 命题在此显然成立.
再利用“相伴层”函子,
按 \hyperrefext[II][II.sec4]{4} 的标准方法过渡到一般情形.

\begin{enonce*}{命题 9.3.4}\etiquette[9.3.4]{IV.prop9.3.4}
沿用 \hyperref[IV.subsubsec9.3.2]{9.3.2} 的记号,
由 $E$ 中满足 $j^{\ast}(X)$ 是 $\mathcal{U}$ 的终对象的对象 $X$
组成的 $E$ 的严格全子范畴 $\mathcal{F}$
\textup{(\cf \hyperref[IV.prop9.3.3]{9.3.3} a))}
是 $E$ 的一个子拓扑斯:
\ie (\hyperref[IV.def9.1.1]{9.1.1}),
它是一个拓扑斯, 且包含函子 $i_{\ast}:\mathcal{F}\rightarrow E$
是某个拓扑斯态射 $i:\mathcal{F}\rightarrow E$ 的正像函子,
其逆像函子为 $X\mapsto i^{\ast}X=X_{\complement U}$.
伴随态射 $X\rightarrow i_{\ast}i^{\ast}X$,
就是 \eqhyperref[IV.eq9.3.2.3]{9.3.2.3} 中的典范态射 $X\rightarrow X_{\complement U}$.
\end{enonce*}

\noindent \emph{证明}:
对 $E$ 的每个对象 $X$,
用 $u(X)$ 表示典范态射 $X\rightarrow X_{\complement U}$.
态射 $u(X)$ 关于 $X$ 具有函子性,
且对每个对象 $X$, $u(X_{\complement U})$ 是同构.
由这两个性质, 形式地、直接地可知,
函子 $i^{\ast}$ 是函子 $i_{\ast}$ 的左伴随,
并且伴随态射为 $u(X)$.
由于函子 $X\rightarrow X_{\complement U}$ 左正合,
函子 $i^{\ast}$ 左正合.
于是由 \hyperrefext[II][II.theo5.5]{5.5} 可知 $\mathcal{F}$ 是拓扑斯,
并由定义 \hyperref[IV.def3.1]{3.1} 可知二元组 $(i^{\ast},i_{\ast})$ 是拓扑斯态射.

\setcounter{subsubsection}{4}
\subsubsection{}\etiquette[9.3.5]{IV.subsubsec9.3.5}
\hyperref[IV.prop9.3.4]{9.3.4} 中描述的 $E$ 的子拓扑斯 $\mathcal{F}$,
称为\emph{$E$ 中与 $E$ 的开子对象 $U$ 互补的闭子拓扑斯},
或等价地, 与 $E$ 中对应于 $U$ 的开子拓扑斯互补的闭子拓扑斯.
按照 \hyperref[IV.def9.1.1]{9.1.1} b),
\hyperref[IV.prop9.3.4]{9.3.4} 中构造的态射
$i:\mathcal{F}\rightarrow E$ 称为\emph{包含态射}.
若对 $E$ 的全子范畴 $\mathcal{F}$,
存在 $E$ 的开子对象 $U$, 使 $\mathcal{F}$ 是与 $U$ 互补的闭子拓扑斯,
就称它是 $E$ 的一个\emph{闭子拓扑斯}.
注意, 这个 $U$ 唯一确定为 $i_{\ast}(\phi_{\mathcal{F}})$,
其中 $\phi_{\mathcal{F}}$ 是 $\mathcal{F}$ 的始对象, 等于 $i^{\ast}(\phi_E)$;
它也称为\emph{$E$ 中与闭子拓扑斯 $\mathcal{F}$ 互补的开子对象},
并且\pageoriginale 由 $U$ 定义的 $E$ 的子拓扑斯,
也称为 \emph{$E$ 中与 $\mathcal{F}$ 互补的开子拓扑斯}.
最后, 若拓扑斯嵌入 (\hyperref[IV.def9.1.1]{9.1.1} a)) $i:F\rightarrow E$
所对应的 $E$ 的子拓扑斯为闭的, 就称它是\emph{闭嵌入}.

若 $G$ 是 $E$ 上的一个群对象, $s\in\Hom(e_E,G)$,
则 $G$ (\resp $s$) 的\emph{支集},
是 $E$ 中与 $G$ (\resp $s$) 的余支集开子对象互补的闭子拓扑斯
(\hyperref[IV.subsubsec8.5.1]{8.5.1}, \hyperref[IV.subsubsec8.5.2]{8.5.2}).

\subsubsection{}\etiquette[9.3.6]{IV.subsubsec9.3.6}
沿用 \hyperref[IV.subsubsec9.3.2]{9.3.2}
和 \hyperref[IV.prop9.3.4]{9.3.4} 的记号,
于是有两个拓扑斯态射
\begin{equation*}
\mathcal{U}\xrightarrow{j}
E\xleftarrow{i}\mathcal{F}{\quoi},\tag{9.3.6.1}
\end{equation*}\etiquette[9.3.6.1]{IV.eq9.3.6.1}
它们定义的五个函子组成两个伴随函子序列,
$(j_{!},j^{\ast},j_{\ast})$ 和 $(i^{\ast},i_{\ast})$:
\[
\xymatrix@C=2cm@R=1.5cm@M=.2cm{ \Uu \ar@<6pt>[r]^{j_!}
\ar@<-13pt>[r]|{j_{\ast}} & E \ar@<3pt>[r]^{i^{\ast}}
\ar@<4pt>[l]|{j^{\ast}}& \Ff
\ar@<13pt>[l]|{i_{\ast}}.}\leqno{(9.3.6.2)}
\]\etiquette[9.3.6.2]{IV.eq9.3.6.2}

函子
\begin{equation*}
\rho=i^{\ast}j_{\ast}:\mathcal{U}\longrightarrow
\mathcal{F}\tag{9.3.6.3}
\end{equation*}\etiquette[9.3.6.3]{IV.eq9.3.6.3}
称为相对于 $E$ 的开子对象 $U$ 的\emph{粘合函子}.
更一般地, 若 $\mathcal{U}$ 和 $\mathcal{F}$ 是任意拓扑斯,
函子 $\rho:\mathcal{U}\rightarrow \mathcal{F}$
若左正合且可及 (\hyperref[I.def9.2]{I.9.2}), 就称为\emph{粘合函子};
或等价地 (\sisi{\hyperref[IV.rem1.8]{1.8}\quad}{\hyperrefext[I][I.sec8]{8}}),
若它有一个 pro-伴随函子, 就称为粘合函子.
这一术语的理由将在下文的 \hyperref[IV.prop9.4.1]{9.4.1} d)
和 \hyperref[IV.theo9.5.4]{9.5.4} 中显现.

% END INCLUDED SGA4_English_current_buildable_source_20260906/source/04/drafts/SGA4_IV_LINES_6402_6597_English.texfrag


% BEGIN INCLUDED SGA4_English_current_buildable_source_20260906/source/04/drafts/SGA4_IV_LINES_6598_6787_English.texfrag
\subsection{闭子拓扑斯 $\mathcal{F}$ 及 $i:\mathcal{F}\rightarrow E$ 的初步性质}\etiquette[9.4]{IV.subsec9.4}

\begin{enonce*}{命题 9.4.1}\etiquette[9.4.1]{IV.prop9.4.1}
记号与 \hyperref[IV.subsubsec9.3.2]{9.3.2}
和 \hyperref[IV.prop9.3.4]{9.3.4} 相同.
\begin{enumerate}
\renewcommand{\theenumi}{\alph{enumi}}
\renewcommand{\labelenumi}{\theenumi)}
\item 函子 $i_{\ast}:\mathcal{F}\rightarrow E$
将满态射族映为满态射族, 并与纤维余积和滤过归纳极限的形成交换.

\item 对 $E$ 中使投射态射 $U\times X\rightarrow U$ 为满态射的每个对象 $X$,
伴随态射 $X\rightarrow i_{\ast}i^{\ast}X$ 是满态射.

\hypertarget{sga4.exp.IV.proposition.9.4.1.item.c}{}\item 从 $E$ 分别到 $\mathcal{F}$
和 $\mathcal{U}$ 的这对\pageoriginale 函子 $(i^{\ast},j^{\ast})$
保守 {\rm(I \hyperref[I.def6.1]{6.1})},
\ie $E$ 的态射 $m$ 是同构,
当且仅当 $i^{\ast}(m)$ 和 $j^{\ast}(m)$ 都是同构
(等价地, 这对函子 $(i^{\ast},j^{\ast})$ 忠实
(\hyperref[IV.subsubsec6.4.0]{6.4.0})).

\item 粘合函子 $\rho=i^{\ast}j_{\ast}:\mathcal{U}\rightarrow\mathcal{F}$
左正合且可及 {\rm(\hyperrefext[I][I.def9.2]{9.2})},
或等价地 {\rm(\hyperrefext[I][I.sec8]{8})},
它有一个 pro-伴随 $\sigma:\Pro\mathcal{F}\rightarrow\Pro\mathcal{U}$.
\end{enumerate}
\end{enonce*}

\noindent \emph{证明}:
a) 由 \hyperref[IV.prop9.3.3]{9.3.3} c) 得出.
对 b), 只须使用 $X_{\complement U}$ 的定义
\eqhyperref[IV.eq9.3.2.2]{9.3.2.2}, 并注意
$i_{\ast}i^{\ast}X=X_{\complement U}$
(\hyperref[IV.prop9.3.4]{9.3.4}).
断言 c) 来自 \hyperref[IV.prop9.3.3]{9.3.3} b).
函子 $i^{\ast}$ 和 $j_{\ast}$ 左正合, 所以 $\rho$ 也左正合.
函子 $j_{\ast}$ 是拓扑斯态射的正像函子, 因而可及
(\hyperrefext[I][I.prop9.5]{9.5});
注意拓扑斯是一个可及范畴
(\hyperrefext[I][I.rem9.11.3]{9.11.3}).
由于函子 $i^{\ast}$ 与归纳极限交换 (它是一个逆像函子),
函子 $\rho$ 作为两个可及函子的复合, 是可及的.

\begin{enonce*}{命题 9.4.2}\etiquette[9.4.2]{IV.prop9.4.2}
设 $E'$ 是一个拓扑斯. 则函子 $h\mapsto i\circ h$,
$$
\Hhomtop(E',\mathcal{F})\longrightarrow \Hhomtop(E',E),
$$
全忠实, 且其本质像由满足 $g^{\ast}(U)$ 为 $E'$ 的始对象 $\phi_{E'}$
的拓扑斯态射 $g:E'\rightarrow E$ 组成.
\end{enonce*}

由 \hyperref[IV.prop9.1.4]{9.1.4}, 所考虑的函子全忠实,
并且 $g$ 属于其本质像,
当且仅当对 $E'$ 的每个对象 $X'$,
$E$ 的对象 $g_{\ast}(X')$ 属于 $\mathcal{F}$,
\ie
\begin{equation*}
j^{\ast}(g_{\ast}(X'))\simeq e_{\mathcal{U}}\text{ 对每个 }X'\in
\Ob E'{\quoi}.\tag{*}
\end{equation*}
置 $U'=g^{\ast}(U)$, 并考虑拓扑斯图
\[
\xymatrix@C=1.5cm{
E'_{/U'}\ar[r]^{j'} \ar[d]_{g_U}& E'\ar[d]^{g}\\
E_{/U}   \ar[r]^{j}             & E\quad , }
\]
显然\pageoriginale 有一个典范同构
$$
j^{\ast}(g_{\ast}(X'))\simeq g_{U\ast}({j'}^{\ast}(X')){\quoi}.
$$
函子 ${j'}^{\ast}$ 也显然本质满射, 所以条件 (*) 等价于
\begin{equation*}
g_{U\ast}(Y')\simeq e_{\mathcal{U}}\text{ 对每个 }Y'\in\Ob
E'_{/U'}{\quoi}.\tag{**}
\end{equation*}
利用 $g_{U\ast}$ 和 $g_U^{\ast}$ 的伴随关系,
立即可见这意味着:
对 $E_{/U}$ 的每个对象 $Y$, $g_U^{\ast}(Y)$ 是 $E'_{/U'}$ 的始对象;
或者利用 $E'_{/U'}$ 的始对象严格这一事实
(\hyperrefext[II][II.def4.5]{4.5}),
只须对 $\mathcal{U}$ 的终对象 $Y=e_{\mathcal{U}}$ 有此性质.
但这也就是说, $E'_{/U'}$ 的终对象 $U'$ 是始对象,
显然这等价于 $U'$ 是 $E'$ 的始对象 $\phi_{E'}$, \cqfd

\begin{enonce*}{推论 9.4.3}\etiquette[9.4.3]{IV.cor9.4.3}
设 $g:E'\rightarrow E$ 是一个拓扑斯态射,
$U'=g^{\ast}(U)$, $\mathcal{U}'=E'_{/U'}$,
$\mathcal{F}'$ 是 $E'$ 中与开子对象 $U'$ 互补的闭子拓扑斯;
最后设 $i':\mathcal{F}'\rightarrow E'$ 为包含态射.
则拓扑斯态射 $g\circ i':\mathcal{F}'\rightarrow E$
在相差同构的意义下经过一个拓扑斯态射
$g_{\mathcal{F}}:\mathcal{F}'\rightarrow\mathcal{F}$ 分解,
且相应拓扑斯图
\[\vcenter{
\xymatrix@=.5cm{ \Ff' \ar[rr]^{i'}\ar[dd]_{g_{\Ff}}  &&
E'\ar[dd]^{g}
\\
                          &            &  \\
\Ff  \ar[rr]^(.4){i}  &  \ar@/^.5pc/[ru]  &   E }}\leqno{(9.4.3.1)}
\]\etiquette[9.4.3.1]{IV.eq9.4.3.1}
是 $2$-Cartesian 的 {\rm(\cf \hyperref[IV.prop5.11]{5.11}).}
\end{enonce*}

这形式地由 \hyperref[IV.prop9.4.2]{9.4.2} 和定义得出.

因此, 用 \hyperref[IV.exer9.1.6]{9.1.6} d) 中引入的术语,
尤其可以说: $E$ 的\emph{闭}子拓扑斯 $\mathcal{F}$
在拓扑斯态射 $g:E'\rightarrow E$ 下的逆像,
是 $E'$ 的一个\emph{闭}子拓扑斯 $\mathcal{F}'$
(即与开子对象 $U'=g^{\ast}(U)$ 互补的闭子拓扑斯,
其中 $U$ 是 $E$ 中与 $\mathcal{F}$ 互补的开子对象).

\begin{enonce*}{推论 9.4.4}\etiquette[9.4.4]{IV.cor9.4.4}
设 $E$ 是一个拓扑斯, $\mathcal{U}$ 是 $E$ 的一个开子拓扑斯,
$\mathcal{F}$ 是互补的闭子拓扑斯.
\begin{enumerate}
\renewcommand{\theenumi}{\alph{enumi}}
\renewcommand{\labelenumi}{\theenumi)}
\item 有\pageoriginale
$$
\mathcal{U}\cap\mathcal{F}\approxeq\Top(∅)\quav
\mathcal{U}\vee\mathcal{F}=E{\quoi},
$$
其中 $\vee$ 表示 $E$ 的所有子拓扑斯所成集合中的 $\Sup$
{\rm(\hyperref[IV.rem9.1.2]{9.1.2} c))}.
(用 \hyperref[IV.exer9.1.13]{9.1.13} 中引入的术语,
$\mathcal{U}$ 和 $\mathcal{F}$ 是 $E$ 的互补子拓扑斯.)

\item $\mathcal{U}$ (\resp $\mathcal{F}$)
是 $E$ 中满足 $E'\cap\mathcal{F}\approxeq\Top(∅)$
(\resp $\mathcal{U}\cap E'\approxeq\Top(∅)$)
的子拓扑斯 $E'$ 中最大的一个.
\end{enumerate}
\end{enonce*}

\noindent \emph{证明.}
\begin{enumerate}
\renewcommand{\theenumi}{\alph{enumi}}
\renewcommand{\labelenumi}{\theenumi)}
\item 第一个关系就是说:
若 $E'$ 是拓扑斯,
且拓扑斯态射 $g:E'\rightarrow E$ 在相差同构的意义下
既经过 $\mathcal{U}$ 又经过 $\mathcal{F}$ 分解,
则 $E'\approxeq\Top(∅)$.
由 \hyperref[IV.subsubsec9.2.5]{9.2.5} 和 \hyperref[IV.prop9.4.2]{9.4.2},
这个假设就是说 $g^{\ast}(U)$ 同时是 $E'$ 的始对象和终对象,
这就证明了断言.
对于第二个关系, 注意由 \hyperref[IV.exer9.1.7.2]{9.1.7.2} e),
它等价于 \hyperref[IV.prop9.4.1]{9.4.1} c).

\item 假设 $E'\cap\mathcal{F}\approxeq\Top(∅)$.
由 \hyperref[IV.cor9.4.3]{9.4.3},
$E'\cap\mathcal{F}$ 等价于 $E'$ 中与 $g^{\ast}(U)$ 互补的子拓扑斯,
其中 $g:E'\rightarrow E$ 是包含态射.
因此该条件就是说 $U'=g^{\ast}(U)$ 是 $E'$ 的终对象,
\ie (\hyperref[IV.subsubsec9.2.5]{9.2.5}) $E'\subset\mathcal{U}$.
假设 $\mathcal{U}\cap E'\approxeq\Top(∅)$.
由 \hyperref[IV.prop5.11]{5.11}, $\mathcal{U}\cap E'$ 等价于 $E'_{/U'}$,
所以该条件就是说 $U'$ 是 $E'$ 的始对象,
\ie (\hyperref[IV.prop9.4.2]{9.4.2}) $E'\subset\mathcal{F}$.
\end{enumerate}

\begin{enonce*}{推论 9.4.5}\etiquette[9.4.5]{IV.cor9.4.5}
设 $(\mathcal{F}_i)_{i\in I}$ 是拓扑斯 $E$ 的一个闭子拓扑斯族,
对每个 $i\in I$, 设 $U_i$ 是 $E$ 中与 $\mathcal{F}_i$ 互补的开子对象.
则将 $\mathcal{F}_i$ 视为 $E$ 的子拓扑斯时所取的交
(\hyperref[IV.prop9.1.3]{9.1.3}) 是一个闭子拓扑斯,
其互补开子对象为 $U=\Sup_iU_i$.
\end{enonce*}

这形式地由交作为 2-积的泛性质和 \hyperref[IV.prop9.4.2]{9.4.2} 得出,
并用到以下事实:
对拓扑斯态射 $g:E'\rightarrow E$,
$g^{\ast}(U)=\Sup_i g^{\ast}(U_i)$ 为始对象,
当且仅当所有 $g^{\ast}(U_i)$ 都为始对象.

类似地可证明:

\begin{enonce*}{推论 9.4.6}\etiquette[9.4.6]{IV.cor9.4.6}
沿用 \hyperref[IV.cor9.4.5]{9.4.5} 的记号,
假设 $I$ 有限.
则 $E$ 的子拓扑斯 $\mathcal{F}_i$ 的上确界\pageoriginale
{\rm(\hyperref[IV.rem9.1.2]{9.1.2} c))}
是 $E$ 的一个闭子拓扑斯, 其互补开子对象为 $\bigcap_iU_i$.
\end{enonce*}


% END INCLUDED SGA4_English_current_buildable_source_20260906/source/04/drafts/SGA4_IV_LINES_6598_6787_English.texfrag


% BEGIN INCLUDED SGA4_English_current_buildable_source_20260906/source/04/drafts/SGA4_IV_LINES_6788_7070_English.texfrag
\begin{enonce*}[remark]{习题 9.4.7}\etiquette[9.4.7]{IV.exer9.4.7}
(\emph{子拓扑斯的粘合}). 设 $E$ 是一个拓扑斯.
\begin{enumerate}
\renewcommand{\theenumi}{\alph{enumi}}
\renewcommand{\labelenumi}{\theenumi)}
\item 设 $E'$ 是 $E$ 的一个子拓扑斯,
$(S_i)_{i\in I}$ 是 $E$ 中一个覆盖终对象的对象族;
对每个 $i\in I$, 设 $S'_i$ 是 $S_i$ 在 $E'$ 中的逆像,
并设 $E'_i\approxeq E'_{/S'_i}$ 是
$E_i=E_{/S_i}$ 中由 $E$ 的子拓扑斯 $E'$
在局部化态射 $E_i\rightarrow E$ 下取逆像而得的子拓扑斯
(\hyperref[IV.exer9.1.6]{9.1.6} d) 和 \hyperref[IV.prop5.11]{5.11}).
证明对 $E$ 的对象 $X$, $X$ 属于 $E'$,
当且仅当对每个 $i\in I$, $X$ 在 $E_i$ 中的逆像 $X| S_i$ 属于 $E'_i$.

\item 假设 $E$ 形如 $C^{\sim}$, 其中 $C$ 是一个 $𝒰$-景.
证明 $C$ 上的预层
\begin{center}
$S\mapsto$ 拓扑斯 $(C_{/S})^{\sim}\approxeq E_{/\varepsilon(S)}$ 的子拓扑斯集合
\end{center}
是一个层. 对任意的 $E$, 由此得出:
存在 $E$ 的一个对象, 代表函子
$$
\text{$S\mapsto$ 拓扑斯 $E_{/S}$ 的子拓扑斯集合;}
$$

\item 沿用 a) 的记号,
证明 $E$ 的子拓扑斯 $E'$ 是开子拓扑斯 (\resp 闭子拓扑斯),
当且仅当对每个 $i\in I$,
$E_i$ 的诱导子拓扑斯 $E'_i$ 也如此.
\end{enumerate}
\end{enonce*}

\begin{enonce*}[remark]{习题 9.4.8}\etiquette[9.4.8]{IV.exer9.4.8}
(\emph{子拓扑斯的内部、外部、闭包和边界.})
设 $E$ 是一个拓扑斯, $E'$ 是一个子拓扑斯,
$i:E'\rightarrow E$ 是包含态射.
\begin{enumerate}
\renewcommand{\theenumi}{\alph{enumi}}
\renewcommand{\labelenumi}{\theenumi)}
\item 证明 $i_{\ast}(\phi_{E'})$ 是 $E$ 中使
$E'| \mathcal{U} \approxeq \Top(∅)$ 的最大的开子对象 $U$,
其中 $\mathcal{U}$ 是 $E$ 中由 $U$ 定义的开子拓扑斯.
这个 $U$ 或 $\mathcal{U}$ 称为 $E$ 的子拓扑斯 $E'$ 的\emph{外部},
$E$ 中与 $U$ 互补的闭子拓扑斯称为 $E'$ 的\emph{闭包}.

\item 证明 $E$ 中存在一个最大的开子对象 $V$,
使 $E$ 的相应开子拓扑斯 $\mathcal{V}$ 包含在 $E'$ 中.
(利用 \hyperref[IV.exer9.2.7]{9.2.7}.)
这个 $V$ 或 $\mathcal{V}$ 称为 $E$ 的子拓扑斯 $E'$ 的\emph{内部}.
$E$ 的开子对象 $U$ 和 $V$ 不交 $(U\cap V=\phi_E)$;
$E$ 中与 $\Sup(U,V)=U\amalg V$ 互补\pageoriginale 的闭子拓扑斯,
称为 $E$ 的子拓扑斯 $E'$ 的\emph{边界}.

\item 子拓扑斯 $E'$ 是 $E$ 的开子拓扑斯 (\resp 闭子拓扑斯),
当且仅当它等于其内部 (\resp 等于其闭包).
$E'$ 的边界为空,
当且仅当 $E'$ 同时是 $E$ 的开子拓扑斯和闭子拓扑斯;
或等价地, 它是 $E$ 中由 $E$ 的某个作为终对象 $e$ 的\emph{直和项}的开子对象 $V$
定义的开子拓扑斯,
\ie 存在 $e$ 的一个子对象 $U$, 使 $e\simeq U\amalg V$.
$E'$ 的外部“为空”, \ie 闭包等于 $E$,
当且仅当 $i:E'\rightarrow E$ 支配 (\hyperref[IV.exer8.8]{8.8}).

\item 设 $S$ 是 $E$ 的一个对象, $F=E_{/S}$ 是诱导拓扑斯,
令 $S'=i^{\ast}(S)$; 则 $F'\simeq E'_{/S'}$ 是由 $E$ 的子拓扑斯 $E'$ 诱导的 $F$ 的子拓扑斯.
证明 $F'$ 的外部 (\resp 闭包, \resp 内部, \resp 边界)
由 $E'$ 的外部 (\resp\ldots) 诱导.
\end{enumerate}
\end{enonce*}

\begin{enonce*}[remark]{习题 9.4.9}\etiquette[9.4.9]{IV.exer9.4.9}
(\emph{拓扑斯的局部闭子拓扑斯.})
若拓扑斯 $E$ 的子拓扑斯 $F$ 是一个开子拓扑斯与一个闭子拓扑斯的交,
则称它是 $E$ 的\emph{局部闭子拓扑斯}.
\begin{enumerate}
\renewcommand{\theenumi}{\alph{enumi}}
\renewcommand{\labelenumi}{\theenumi)}
\item 证明 $E$ 的有限族局部闭子拓扑斯的交是局部闭的.
若 $g:E'\rightarrow E$ 是一个拓扑斯态射,
证明 $E$ 的局部闭子拓扑斯的逆像
(\hyperref[IV.exer9.1.6]{9.1.6} d))
是 $E'$ 的局部闭子拓扑斯.

\item 设 $X$ 是一个拓扑空间, $E=\Top(X)$.
对 $X$ 的每个子集 $X'$, 考虑 $E$ 的子拓扑斯 $T(X')$,
它是 $\Top(X')$ 在 $\Top(X)$ 中的本质像.
证明若 $X'$ 是局部闭子集, 则 $T(X')$ 是 $E$ 的局部闭子拓扑斯;
若再假设 $X$ 和 $X'$ 清醒, 逆命题也成立.
证明 $X'\mapsto T(X')$ 在 $X$ 的局部闭子集 $X'$ 所成的集合,
与 $E$ 的局部闭子拓扑斯所成的集合之间定义一个序同构.

\item 设 $F$ 是 $E$ 的一个子拓扑斯.
证明 $F$ 的每个开子拓扑斯\pageoriginale
都由 $E$ 的一个开子拓扑斯诱导.
(利用 \hyperref[IV.exer9.1.8]{9.1.8} c).)
证明 $F$ 是 $E$ 的局部闭子拓扑斯,
当且仅当它是其闭包的开子拓扑斯 (\hyperref[IV.exer9.4.8]{9.4.8} a)).
再利用 \hyperref[IV.exer9.4.8]{9.4.8} d) 得出:
$F$ 是 $E$ 的局部闭子拓扑斯的性质是 $E$ 上的局部性质.

\item 设 $F$ 是 $E$ 的一个局部闭子拓扑斯.
证明在将 $F$ 写成 $E$ 的开子拓扑斯 $\mathcal{U}$
与闭子拓扑斯 $\mathcal{F}$ 的交的所有方式中,
存在一种使 $\mathcal{U}$ 尽可能大而 $\mathcal{F}$ 尽可能小.
(取 $\mathcal{F} =$ $F$ 的闭包,
$\mathcal{U} =$ 在 $\mathcal{F}$ 上诱导子拓扑斯 $F$ 的最大开子拓扑斯.)
证明 $\mathcal{U}$ 和 $\mathcal{F}$ 的形成与 $E$ 上的局部化相容.
\end{enumerate}
\end{enonce*}

\begin{enonce*}[remark]{习题 9.4.10}\etiquette[9.4.10]{IV.exer9.4.10}
仿照拓扑空间中熟悉的相应概念
(EGA ${\rm 0_{III}}$ 9.1, EGA IV 10.1),
建立\emph{可构造、局部可构造、拟可构造和局部拟可构造子拓扑斯}的概念.
对于形如 $\Top(X)$ 的拓扑斯,
可重新得到 $X$ 的可构造子集 (\resp\ldots) 所成的集合,
与 $\Top(X)$ 的可构造子拓扑斯 (\resp\ldots) 所成的集合之间的一个双射.
\end{enonce*}

\subsection{粘合定理}\etiquette[9.5]{IV.subsec9.5}

\subsubsection{}\etiquette[9.5.1]{IV.subsubsec9.5.1}
设 $f:A\rightarrow B$ 是两个范畴之间的一个函子.
用 $(B,A,f)$ 表示如下范畴:
$(B,A,f)$ 的对象是三元组
$$
(X,Y,u)
$$
其中 $X$ 是 $B$ 的对象, $Y$ 是 $A$ 的对象,
$u$ 是从 $X$ 到 $f(Y)$ 的态射;
两个对象 $(X,Y,u)$ 和 $(X',Y',u')$ 之间的态射是二元组 $(m,m')$,
其中 $m$ 是从 $X$ 到 $X'$ 的态射,
$m'$ 是从 $Y$ 到 $Y'$ 的态射, 并使下图交换:
\[
\xymatrix@C=1.5cm@R=1cm{
X \ar[r]^{u} \ar[d]_{m} & **[r]f(Y) \ar[d]^{f(m')}\\
X' \ar[r]^{u'}& **[r]f(Y') }
\]
态射的\pageoriginale 复合按显然的方式定义.

\subsubsection{}\etiquette[9.5.2]{IV.subsubsec9.5.2}
注意, 范畴 $(B,A,f)$ 函子性地依赖于数据系 $A$, $B$, $f$;
其准确表述留给读者.
特别地, 若 $f'$ 是从 $A$ 到 $B$ 的另一个函子,
则每个函子态射
$$
f\xrightarrow{u} f'
$$
定义一个函子 $(B,A,f)\rightarrow(B,A,f')$;
若 $u:f\rightarrow f'$ 是同构, 这个函子就是同构.

\subsubsection{}\etiquette[9.5.3]{IV.subsubsec9.5.3}
回到 \hyperref[IV.subsec9.3]{9.3} 的情形,
考虑 \hyperref[IV.subsubsec9.5.1]{9.5.1} 中定义的范畴
$(\mathcal{F},\mathcal{U},\rho)$.
与 $E$ 的每个对象 $X$ 相应的,
有 $(\mathcal{F},\mathcal{U},\rho)$ 的一个对象
$(i^{\ast}X,j^{\ast}X,h(X):i^{\ast}X\rightarrow\rho j^{\ast}X)$,
其中 $h(X)$ 由对伴随态射 $X\rightarrow j_{\ast}j^{\ast}(X)$
应用函子 $i^{\ast}$ 得到
(有 $\rho=i^{\ast}j_{\ast}$
\sisi{(9.10.)}{\hyperref[IV.eq9.3.6.3]{9.3.6.3}}).
由此定义一个函子
\begin{equation*}
\Phi:E\longrightarrow (\mathcal{F}, \mathcal{U},
\rho){\quoi}.\tag{9.5.3.1}
\end{equation*}\etiquette[9.5.3.1]{IV.eq9.5.3.1}

\begin{enonce*}{定理 9.5.4}\etiquette[9.5.4]{IV.theo9.5.4}
\begin{enumerate}
\renewcommand{\theenumi}{\alph{enumi}}
\renewcommand{\labelenumi}{\theenumi)}
\item 设 $E$ 是一个拓扑斯, $\mathcal{U}$ 是 $E$ 的一个开子拓扑斯,
$\mathcal{F}$ 是互补的闭子拓扑斯.
函子
$$
\rho=i^{\ast}j_{\ast}:\mathcal{U} \longrightarrow \mathcal{F}
$$
左正合且可及,
并且函子 $\Phi$ $\hyperref[IV.eq9.5.3.1]{(9.5.3.1)}$ 是一个范畴等价.

\item 反过来, 设 $\mathcal{U}$ 和 $\mathcal{F}$ 是两个拓扑斯,
$\rho:\mathcal{U}\rightarrow\mathcal{F}$ 是一个粘合函子
$(\hyperref[IV.subsubsec9.3.6]{9.3.6})$
(\ie 可及 {\rm (\hyperrefext[I][I.def9.2]{9.2})} 且左正合).
范畴 $E=(\mathcal{F},\mathcal{U},\rho)$ 是一个拓扑斯.
设 $\phi_{\mathcal{F}}$ 是 $\mathcal{F}$ 的一个始对象,
$e_{\mathcal{U}}$ 是 $\mathcal{U}$ 的一个终对象,
并设 $U=(\phi_{\mathcal{F}},e_{\mathcal{U}},
\phi_{\mathcal{F}}\rightarrow\rho(e_{\mathcal{U}}))\in\Ob E$.
函子 $X\mapsto(\phi_{\mathcal{F}},X,\phi_{\mathcal{F}}\rightarrow\rho(X))$
诱导 $\mathcal{U}$ 与 $E_{/U}$ 的一个等价,
从而与 $E$ 中由 $U$ 定义的开子拓扑斯 $\mathcal{U}'$ 等价.
函子 $Y\mapsto(Y,e,Y\rightarrow\rho(e))$ 是闭嵌入
$(\hyperref[IV.subsubsec9.3.5]{9.3.5})$,
\ie 它诱导 $\mathcal{F}$ 与 $E$ 的某个闭子拓扑斯 $\mathcal{F}'$ 的一个等价.
子拓扑斯 $\mathcal{U}'$ 和 $\mathcal{F}'$\pageoriginale 互补.
设 $\rho':\mathcal{U}'\rightarrow\mathcal{F}'$ 是粘合函子.
函子 $\rho$ 和 $\rho'$ 与上述等价相容.
\end{enumerate}
\end{enonce*}

\begin{enonce*}{定义 9.5.5}\etiquette[9.5.5]{IV.def9.5.5}
设 $\mathcal{U}$ 和 $\mathcal{F}$ 是两个拓扑斯,
$\rho:\mathcal{U}\rightarrow\mathcal{F}$ 是一个粘合函子
$(\hyperref[IV.subsubsec9.3.6]{9.3.6})$.
拓扑斯 $(\mathcal{F},\mathcal{U},\rho)$
称为通过函子 $\rho$ 将 $\mathcal{U}$ 和 $\mathcal{F}$ 粘合而构造的拓扑斯.
\end{enonce*}

\setcounter{subsubsection}{5}
\subsubsection{}\etiquette[9.5.6]{IV.subsubsec9.5.6}
\emph{\hyperref[IV.theo9.5.4]{9.5.4} a) 的证明}
\footnote[1]{关于比 \hyperref[IV.def9.5.5]{9.5.5} 更一般的粘合陈述,
\cf 习题\stepcounter{cnote}\textsuperscript{\arabic{cnote}}%
\textup{ [N.D.E.: 这道习题似乎并不存在.]}.}
第一个断言只是重述 \hyperref[IV.prop9.4.1]{9.4.1} d).
这对函子 $(i^{\ast},j^{\ast})$ 保守,
或等价地忠实 (\hyperref[IV.prop9.4.1]{9.4.1} c)).
因此, 函子 $\Phi$ 更是忠实的.
下面证明它全忠实.
由 \hyperref[IV.prop9.3.3]{9.3.3} b) 可知,
对 $E$ 的每个对象 $X$, 下图为 Cartesian 图:
\[\vcenter{
\xymatrix@C=2.5cm{
X  \ar[r] \ar[d]               & j_{\ast}j^{\ast}X \ar[d]\\
i_{*}i^{*}X   \ar[r]     & i_{\ast}i^{\ast}j_{\ast}j^{\ast}X ;
}}\leqno{(*)}
\]
设 $X$ 和 $Y$ 是 $E$ 的两个对象, 并设 $(m,m')$:
\[\vcenter{
\xymatrix@C=2.75cm@R=1cm{
 i^{\ast}X  \ar[r] \ar[d]_{m}               & i^{\ast}j_{\ast}j^{\ast}X \ar[d]_{i^{\ast}j_{\ast}m'}\\
i^{*}Y   \ar[r]     & i^{\ast}j_{\ast}j^{\ast}Y ; }}
\]
是从 $\Phi(X)$ 到 $\Phi(Y)$ 的态射.
利用纤维积的函子性以及图 $(\ast)$,
得到一个使下列图交换的态射 $w:X\rightarrow Y$:
\begin{displaymath}
\xymatrix@C=.5cm@R=1.25cm{
X \ar@{-}[rr]^{w} \ar[d]  &             &Y \ar[d]\\
 j_{\ast}j^{\ast}X \ar[rr]^{j_{\ast}m'}       &   &
j_{\ast}j^{\ast}Y , }\qquad\qquad \xymatrix@C=.5cm@R=1.25cm{
X \ar@{-}[rr]^{w} \ar[d]& &Y \ar[d]\\
 i_{\ast}i^{\ast}X \ar[rr]^{i_{\ast}m}&   & i_{\ast}i^{\ast}Y
}
\end{displaymath}
分别\pageoriginale 应用函子 $j^{\ast}$ 和 $i^{\ast}$,
得到 $j^{\ast}w=m'$ 和 $i^{\ast}w=m$.
现在证明函子 $\Phi$ 本质满射.
设 $(Y,X,u:Y\rightarrow i^{\ast}j_{\ast}(X))$
是 $(\mathcal{F},\mathcal{U},i^{\ast}j_{\ast})$ 的一个对象.
由此得到一个图:
\[
\xymatrix@C=2cm{
                             & j_{\ast}X \ar[d]\\
i_{\ast}Y \ar[r]^{i_{\ast}u} & i_{\ast}i^{\ast}j_{\ast}X }
\]
于是取纤维积, 得到 Cartesian 图:
\[\vcenter{
  \xymatrix@C=2.75cm@R=1.25cm{
W \ar[r] \ar[d]& j_{\ast}X \ar[d]\\
i_{\ast}Y \ar[r]^{i_{\ast}u}& i_{\ast}i^{\ast}j_{\ast} X\\
}}\leqno{(**)}
\]

留待读者证明 $\Phi(W)$ 同构于
$(Y,X,u:Y\rightarrow i^{\ast}j_{\ast}(X))$,
且图 (${\ast}{\ast}$) 同构于与 $W$ 相应的图 ($\ast$).

\hypertarget{sga4.exp.IV.subsubsection.9.5.7}{}\subsubsection{}

% END INCLUDED SGA4_English_current_buildable_source_20260906/source/04/drafts/SGA4_IV_LINES_6788_7070_English.texfrag


% BEGIN INCLUDED SGA4_English_current_buildable_source_20260906/source/04/drafts/SGA4_IV_LINES_7071_7243_English.texfrag
\emph{\hyperref[IV.theo9.5.4]{9.5.4} \sisi{1)}{b)} 的证明}.
我们只证明 $(\mathcal{F},\mathcal{U},\rho)$ 是一个拓扑斯.
其他断言的证明留给读者.
验证 \hyperref[IV.subsubsec1.1.2]{1.1.2} 的性质 a), b), c), d).
\begin{enumerate}
\renewcommand{\theenumi}{\alph{enumi}}
\renewcommand{\labelenumi}{\theenumi)}
\item \emph{有限投射极限可表}:
这很清楚, 因为范畴 $\mathcal{U}$ 和 $\mathcal{F}$ 是拓扑斯,
且 $\rho$ 与有限投射极限交换.

\item \emph{直和可表、不交且普遍}:
设 $(Y_i\xrightarrow{u_i}\rho(X_i))$, $i\in I$,
是 $(\mathcal{F},\mathcal{U},\rho)$ 的一个对象族.
根据归纳极限的定义, 有一个典范态射
$$
\coprod_{i\in I}\rho(X_i)\longrightarrow\rho(\coprod_{i\in I}X_i),
$$
因而,\pageoriginale 与 $\coprod u_i$ 复合, 得到一个态射
$$
n:\coprod_{i\in I}Y_i\longrightarrow\rho(\coprod_{i\in I}X_i){\quoi}.
$$
直接可验证, 最后这个对象是所给族的直和,
且这个直和不交并且普遍.

\item \emph{等价关系有效且普遍}:

设
\[
\xymatrix@C=2cm@R=1.5cm{
R_1 \ar@<3pt>[r] \ar@<-3pt>[r] \ar[d]& Y \ar[d]^{w}\\
\rho(R_2) \ar@<3pt>[r]^{\rho(u)} \ar@<-3pt>[r]_{\rho(v)} & \rho(X) }
\]
是对象 $Y\xrightarrow{w}\rho(X)$ 上的一个等价关系.
则 $\xymatrix@C=.6cm{R_1\ar@<-2pt>[r]\ar@<2pt>[r]&Y}$
是 $\mathcal{F}$ 中 $Y$ 上的一个等价关系,
而 $\xymatrix{R_2\ar@<2pt>[r]^{u}\ar@<-2pt>[r]_{v}&X\\}$
是 $\mathcal{U}$ 中 $X$ 上的一个等价关系.
商 $Y/R_1$, $\rho(X)/\rho(R_2)$ 和 $X/R_2$
在 $\mathcal{F}$ 和 $\mathcal{U}$ 中可表,
因为这些范畴是拓扑斯.
此外, 根据归纳极限的定义,
典范态射 $\rho(X)\rightarrow\rho(X/R_2)$
经过 $\rho(X)/\rho(R_2)$ 分解.
由归纳极限的函子性, 得到一个典范态射
$Y/R_1\xrightarrow{t}\rho(X/R_2)$.
可验证, 最后这个对象是 $Y\xrightarrow{w}\rho(X)$ 对所考虑的等价关系的商,
这个商有效 (I~10), 且所有这些性质在基变换下保持.

\item $(\mathcal{F},\mathcal{U},\rho)$ \emph{有一个小生成族}.
这个性质由 \sisi{(1.9.2 5)}{\hyperrefext[I][I.cor9.25]{9.25}} 得出
(注意 $\rho$ 可及), 只须取
$B=\Delta_1=(0\xrightarrow{f}1)$, $E_0=\mathcal{F}$,
$E_1=\mathcal{U}$ 及 $f^{\ast}=\rho:E_1\rightarrow E_0$,
于是 $\Hhom_B(B,E)=(\mathcal{F},\mathcal{U},\rho)$.

也可避免使用 \sisi{(1.9.2 5)}{\hyperrefext[I][I.cor9.25]{9.25}}
(前面给出的证明颇为艰难!),
而将关于 $\rho$ 的假设写成 $\rho$ 有一个 pro-伴随的形式:
$$
\sigma:\Pro(\mathcal{F})\longrightarrow\Pro(\mathcal{U}){\quoi},
$$
(这种\pageoriginale 解释 (\hyperrefext[I][I.sec8]{8})
以 \hyperrefext[I][I.sec9]{9} 为基础, 后者的证明明显容易理解得多).
只须注意: 若 $C'$ (\resp $C''$) 是 $\mathcal{U}$
(\resp $\mathcal{F}$) 的一个生成子范畴,
并且对每个 $X''\in\Ob C''$,
将 $\sigma(X'')$ 表示为一个投射系 $(\sigma(X'')_i)_{i\in I(X'')}$,
其中 $I(X'')$ 是一个小滤过有序集,
则由以下两种形式的对象组成的 $E$ 的对象族是一个生成族:
一种形如 $(\phi_{\mathcal{F}},X',\phi_{\mathcal{F}}\rightarrow\rho(X'))$,
其中 $X'\in\Ob C'$;
另一种形如 $(X'',\sigma(X'')_i,u_{X''}:X''\rightarrow\rho(\sigma(X'')_i))$,
其中 $X''\in\Ob C''$, $i\in I(X'')$,
而 $u_{X''}$ 在伴随下对应于典范态射 $\sigma(X'')\rightarrow\sigma(X'')_i$.
当 $C'$ 和 $C''$ 取为小时, 这个生成族显然是小的.
验证确实直接, 留给读者.
\end{enumerate}

\setcounter{subsubsection}{4}\setcounter{paragraph}{7}
\paragraph{}\etiquette[9.5.4.8]{IV.par9.5.4.8}
沿用 \hyperlink{sga4.exp.IV.subsubsection.9.5.7}{9.5.7} b) 的记号,
说明当 $E=(\mathcal{F},\mathcal{U},\rho)$ 时,
如何在相差函子的典范同构的意义下,
把函子系 \eqhyperref[IV.eq9.3.6.2]{9.3.6.2} 显式写出.
下面断言的详细验证
(它们本质上都是 \hyperref[IV.theo9.5.4]{9.5.4} 的机械应用) 留给读者.

\setcounter{subsubsection}{8}\setcounter{paragraph}{0}
\paragraph{}\etiquette[9.5.8.1]{IV.par9.5.8.1}
函子
$$
i^{\ast}:(\mathcal{F},\mathcal{U},\rho)\longrightarrow\mathcal{F}
$$
由
$$
i^{\ast}(X,Y,u:X\longrightarrow\rho(Y))=X{\quoi}.
$$
给出.

\paragraph{}\etiquette[9.5.8.2]{IV.par9.5.8.2}
函子
$$
i_{\ast}:\mathcal{F}\longrightarrow(\mathcal{F},\mathcal{U},\rho)
$$
由 $i_{\ast}(X)=(X,e,X\rightarrow\rho(e))$ 给出,
其中 $e$ 是 $\mathcal{U}$ 的终对象.

\paragraph{}\etiquette[9.5.8.3]{IV.par9.5.8.3}
伴随态射
$$
\id\longrightarrow i_{\ast}i^{\ast}
$$
同构于函子态射
\[
\xymatrix@C=2cm{
X \ar[r]^{u} \ar[d]_{id}& \rho(Y) \ar[d]\\
X \ar[r]                & \rho(e) }
\]

\paragraph{}\etiquette[9.5.8.4]{IV.par9.5.8.4}
函子\pageoriginale
$$
j^{\ast}:(\mathcal{F},\mathcal{U},\rho)\longrightarrow\mathcal{U}
$$
由\qquad $j^{\ast}(X,Y,u:X\longrightarrow\rho(Y))=Y$ 给出.

\paragraph{}\etiquette[9.5.8.5]{IV.par9.5.8.5}
函子
$$
j_{!}:\mathcal{U}\longrightarrow(\mathcal{F},\mathcal{U},\rho)
$$
由\qquad
$j_{!}(Y)=(\phi_{\mathcal{F}},Y,
\phi_{\mathcal{F}}\longrightarrow\rho(Y))$
给出
\noindent (其中 $\phi_{\mathcal{F}}$ 是 $\mathcal{F}$ 的始对象).

\paragraph{}\etiquette[9.5.8.6]{IV.par9.5.8.6}
函子
$$
j_{\ast}:\mathcal{U}\longrightarrow(\mathcal{F},\mathcal{U},\rho)
$$
由\qquad
$j_{\ast}(Y)=(\rho(Y),Y,\id:\rho(Y)\longrightarrow\rho(Y))$ 给出.

\paragraph{}\etiquette[9.5.8.7]{IV.par9.5.8.7}
伴随态射
$$
j_{!}j^{\ast}\longrightarrow\id
$$
同构于函子态射
\[
\xymatrix{
\phi_{\Ff} \ar[r] \ar[d] & \rho(Y) \ar[d]^{\rho(\id)}\\
X \ar[r]^{u}                 & \rho(Y)     &. }
\]

\paragraph{}\etiquette[9.5.8.8]{IV.par9.5.8.8}
伴随态射
$$
\id\longrightarrow j_{\ast}j^{\ast}
$$
同构于函子态射
\[
\xymatrix{
X  \ar[r]^{u} \ar[d]_{u}& \rho(Y) \ar[d]^{\rho(\id)}\\
\rho(Y) \ar[r]^{\id}& \rho(Y) }
\]

% END INCLUDED SGA4_English_current_buildable_source_20260906/source/04/drafts/SGA4_IV_LINES_7071_7243_English.texfrag


% BEGIN INCLUDED SGA4_English_current_buildable_source_20260906/source/04/drafts/SGA4_IV_LINES_7244_7443_English.texfrag
\begin{enonce*}[remark]{习题 9.5.9}\etiquette[9.5.9]{IV.exer9.5.9}
(\emph{粘合拓扑斯的正合性性质.})
设 $\mathcal{U}$, $\mathcal{F}$ 是两个拓扑斯,
$\rho:\mathcal{U}\rightarrow\mathcal{F}$ 是一个粘合函子,
$E=(\mathcal{F},\mathcal{U},\rho)$ 是粘合拓扑斯,
$j:\mathcal{U}\rightarrow E$ 和 $i:\mathcal{F}\rightarrow E$ 是典范拓扑斯态射,
因而 $\rho=i^{\ast}j_{\ast}$.
\begin{enumerate}
\renewcommand{\theenumi}{\alph{enumi}}
\renewcommand{\labelenumi}{\theenumi)}
\item 证明 $i^{\ast}$ 与 (小) 积交换
(\ie (\hyperref[IV.rem1.8]{1.8}
或 \sisi{\hyperref[IV.rem1.8]{1.8}}{\hyperrefext[I][I.sec8]{8}})
函子 $i^{\ast}$ 有一个左伴随 $i_{!}$),
当且仅当 $\rho$ 也如此,
\ie (\loccit) 当且仅当 $\rho$ 有左伴随
$\sigma:\mathcal{F}\rightarrow\mathcal{U}$
(或等价地, 当且仅当 $\rho$ 的 pro-伴随限制\pageoriginale 到 $\mathcal{F}$
\eqhyperref[IV.eq9.6.3]{9.6.3} 后,
在相差同构的意义下经过 $\mathcal{U}$ 分解):
此时有 $\sigma=j^{\ast}i_{!}$.
因此, 具有这个性质的粘合函子 $\rho$,
在相差同构的意义下对应于有右伴随的函子
$\sigma:\mathcal{F}\rightarrow\mathcal{U}$,
\ie (\hyperref[IV.cor1.6]{1.6}) 与 $𝒰$-小归纳极限交换的函子,
或等价地, 连续函子.
当 $\mathcal{F}$ 等价于拓扑斯 $C^{\sim}$,
其中 $C$ 是一个 $𝒰$-景时,
这些函子在相差同构的意义下对应于连续函子 $C\rightarrow\mathcal{U}$
(\hyperrefext[III][III.prop1.2]{1.2} 和 \hyperrefext[III][III.prop1.7]{1.7}).

\item 假设函子 $i_{!}$ 有定义.
证明 $i_{!}$ 与某一指定类型的小投射极限交换,
当且仅当 $\sigma:\mathcal{F}\rightarrow\mathcal{U}$ 与它们交换.
(对于必要性, 利用 $j^{\ast}$ 与小投射极限交换这一事实;
对于充分性, 利用以 $\sigma$ 表述的 $E$ 的等价描述:
它由三元组 $(Y,X,u)$ 组成,
其中 $Y\in\Ob\mathcal{F}$, $X\in\Ob\mathcal{U}$,
$u:\sigma(Y)\rightarrow X$.)

\item 由 a) 和 b) 给出闭拓扑斯嵌入 $i:\mathcal{F}\rightarrow E$ 的一个例子,
使 $i_{!}$ 存在, 却既不与纤维积交换, 也不与滤过投射极限交换.
(只须找到两个拓扑斯之间的一个连续函子,
使它不具有这两种正合性性质中的任何一种.)
\end{enumerate}
\end{enonce*}

\begin{enonce*}[remark]{习题 9.5.10}\etiquette[9.5.10]{IV.exer9.5.10}
(\emph{形如 $\widehat{C}$ 的拓扑斯的粘合.})
设 $C$ 是一个小范畴, $E=\widehat{C}$, 它是一个 $𝒰$-拓扑斯.
\begin{enumerate}
\renewcommand{\theenumi}{\alph{enumi}}
\renewcommand{\labelenumi}{\theenumi)}
\item 回顾一下, $E$ 的开子对象 $U$ 对应于 $C$ 的\emph{筛} $C'$
\sisi{(1.4.1, \hyperref[IV.subsec4.2]{4.2})}{(\hyperrefext[I][I.def4.1]{4.1}, \hyperref[IV.subsubsec8.4.4]{8.4.4})},
此时 $E_{/U}$ 等价于 $\widehat{C'}$,
函子 $j^{\ast}:E\rightarrow E_{/U}$ 可与限制函子
\hyperref[IV.subsec5.4]{5.4} 等同,
\ie 包含态射 $j:E_{/U}\rightarrow E$ 是
$\widehat{f}:\widehat{C'}\rightarrow\widehat{C}$
(\hyperref[IV.subsubsec4.6.1]{4.6.1}),
其中 $f:C'\rightarrow C$ 为包含.
设 $C''$ 是 $C$ 中与 $C'$ 互补的全子范畴
(\ie $\Ob C''$ 是 $\Ob C'$ 在 $\Ob C$ 中的补集),
$g:C''\rightarrow C$ 是包含.
则 $E=\widehat{C}$ 中与由 $C'$ 定义的开子拓扑斯互补的闭子拓扑斯 $\mathcal{F}$,
典范等价于 $\widehat{C''}$,
且包含态射 $i:\mathcal{F}\rightarrow E$
可与 $\widehat{g}:\widehat{C''}\rightarrow\widehat{C}$ 等同.

\item 考虑\pageoriginale 函子
\begin{equation*}
h:{C'}^0\times C''\longrightarrow (\Ens),\quad
h(X',X'')=\Hom_C(X',X''){\quoi},\tag{9.5.10.1}
\end{equation*}\etiquette[9.5.10.1]{IV.eq9.5.10.1}
并注意, 可以由范畴 $C'$, $C''$ 和双函子 $h$
(后者可以任意选取),
在相差同构的意义下重构 $C$ 及其全子范畴 $C'$.
(比较下文 \hyperref[IV.subsec9.7.1]{9.7.1}.)
给定 $h$ 等价于给定一个函子
$C''\rightarrow\widehat{C'}=\Hhom({C'}^0,(\Ens))$,
或等价地 (\hyperrefext[III][III.prop1.2]{1.2} 和 \hyperrefext[III][III.prop1.7]{1.7}),
给定一个连续函子
\begin{equation*}
\sigma:\widehat{C''}\longrightarrow \widehat{C'}\tag{9.5.10.2}
\end{equation*}\etiquette[9.5.10.2]{IV.eq9.5.10.2}
\ie (\hyperref[IV.cor1.6]{1.6}) 与小归纳极限交换的函子,
或等价地, 有一个右伴随
\begin{equation*}
\rho:\widehat{C}'\longrightarrow \widehat{C''}{\quoi}.\tag{9.5.10.3}
\end{equation*}\etiquette[9.5.10.3]{IV.eq9.5.10.3}
证明 $\rho$ 正是与 $E=\widehat{C}$ 的开子拓扑斯 $\widehat{C}'$ 相应的粘合函子.

\item 由 a) 和 b) 得出:
若给定两个预层范畴之间的一个粘合函子
\eqhyperref[IV.eq9.5.10.3]{9.5.10.3}, 从而给定一个粘合拓扑斯 $E$,
则以下条件等价:
\begin{itemize}
\item[(i)] $E$ 等价于形如 $\widehat{C}$ 的拓扑斯.

\item[(ii)] 函子 $\rho$ 与小积交换
(或等价地, 它有左伴随 \eqhyperref[IV.eq9.5.10.2]{9.5.10.2}).

\item[(iii)] 包含态射 $i:\widehat{C''}\rightarrow E$ 是“本质的”,
\ie $i^{\ast}$ 与积交换, 或等价地, $i^{\ast}$ 有左伴随 $i_{!}$.
\end{itemize}
当这些条件成立时,
利用由 \eqhyperref[IV.eq9.5.10.2]{9.5.10.2}
在 $C''$ 上的限制所定义的双函子 \eqhyperref[IV.eq9.5.10.1]{9.5.10.1},
显式描述一个使 $E\approx\widehat{C}$ 的范畴 $C$.

\item 证明由两个点拓扑斯粘合得到的拓扑斯,
未必等价于形如 $\widehat{C}$ 的拓扑斯.
\end{enumerate}
\end{enonce*}

\begin{enonce*}[remark]{习题 9.5.11}\etiquette[9.5.11]{IV.exer9.5.11}
(从粘合拓扑斯到拓扑斯的态射.)
\begin{enumerate}
\renewcommand{\theenumi}{\alph{enumi}}
\renewcommand{\labelenumi}{\theenumi)}
\item 设 $\rho\sisi{}{: \mathcal{U}}\rightarrow\mathcal{F}$
是两个范畴之间的一个函子,
$E=(\mathcal{F},\mathcal{U},\rho)$\pageoriginale
是 \hyperref[IV.subsubsec9.5.1]{9.5.1} 中的“粘合”范畴,
$F$ 是任意范畴.
考虑函子 $f\mapsto\rho\circ f$:
$$
\bar{\rho}:\bar{\mathcal{U}}=\Hhom(F,\mathcal{U})\longrightarrow
\bar{\mathcal{F}}=\Hhom(F,\mathcal{F}){\quoi},
$$
并定义一个范畴等价
$$
\Hhom(F,E)\xrightarrow{\approx}(\bar{\mathcal{F}},\bar{\mathcal{U}},\bar{\rho}){\quoi}.
$$

\item 设 $f:F\rightarrow E$ 是一个函子,
考虑它与典范函子 $j^{\ast}:E\rightarrow\mathcal{U}$
及 $i^{\ast}:E\rightarrow\mathcal{F}$ 的复合.
证明 $f$ 与某一指定类型的归纳极限交换
(假设这些极限在 $\mathcal{U}$ 和 $\mathcal{F}$ 中可表, 从而在 $E$ 中可表),
当且仅当 $j^{\ast}f$ 和 $i^{\ast}f$ 与它们交换.
证明若 $\rho$ 与某一指定类型的投射极限交换
(于是这些极限在 $E$ 中可表),
则 $f$ 与它们交换, 当且仅当 $j^{\ast}f$ 和 $i^{\ast}f$ 也如此.

\item 现在假设 $\mathcal{U}$ 和 $\mathcal{F}$ 是拓扑斯,
且 $\rho$ 是粘合函子.
设 $F$ 是一个拓扑斯, $f:F\rightarrow E$ 是一个函子.
由 b) 得出: $f$ 是与某个拓扑斯态射 $E\rightarrow F$ 相应的逆像函子,
当且仅当函子 $j^{\ast}f:F\rightarrow\mathcal{U}$
和 $i^{\ast}f:F\rightarrow\mathcal{F}$ 也如此.
由此及 a) 得出:
范畴 $\Hhomtop(E,F)$ 与三元组 $(f',f'',u)$ 所成的范畴之间存在一个等价,
其中 $f'$ (\resp $f''$) 是 $\Hhomtop(\mathcal{U},F)$
(\resp $\Hhomtop(\mathcal{F},F)$) 的一个对象,
而 $u:{f''}^{\ast}\rightarrow\rho\circ{f'}^{\ast}$
是函子 $F\rightarrow\mathcal{F}$ 之间的一个同态:
若 $f$ 是 $\Hhomtop(E,F)$ 的一个对象,
$(f',f'',u)$ 是相应的三元组,
则 $f'=f\circ j$, $f''=f\circ i$.

\item 说明在何种意义下,
可以说 c) 通过三元组 $(\mathcal{U},\mathcal{F},\rho)$
给出了拓扑斯 $E$ 在相差等价
(在某个拓扑斯 2-范畴中) 意义下的一个刻画.
\end{enumerate}
\end{enonce*}

\begin{enonce*}[remark]{习题 9.5.12}\etiquette[9.5.12]{IV.exer9.5.12}
(\emph{两个拓扑空间的粘合.})
\begin{enumerate}
\renewcommand{\theenumi}{\alph{enumi}}
\renewcommand{\labelenumi}{\theenumi)}
\item 设 $X$ 是一个拓扑空间,
$\rho:\Top(X)\rightarrow(\Ens)=\Top(pt)$ 是一个粘合函子,
\ie (\hyperrefext[I][I.sec8]{8}) 一个 pro-可表函子;
设 $A\in\Ob\Pro(\Top(X))$ 是 pro-代表它的对象,
$E$ 是从 $\rho$ 粘合得到的拓扑斯.
证明 $E$ 等价于形如 $\Top(Y)$ 的拓扑斯,
当且仅当 $A$ 同构于 $\Pro(\Ouv(X))$ 的一个对象.
由此\pageoriginale 得出:
若 $X$ 非空, 可以找到 $\rho$ 使 $E$ 不等价于形如 $\Top(Y)$ 的拓扑斯.

\item 更一般地, 考虑一个粘合函子
$\rho:\Top(X)\rightarrow\Top(Y)$,
其中 $X$, $Y$ 是两个拓扑空间.
由它导出一个映射
$$
\phi:Y\longrightarrow \Ob \Pro(\Top(X))
$$
(映射 $Y\rightarrow\Ob\Ppoint(\Top(Y))\rightarrow\Ob\Pro(\Top(Y))$
(\hyperref[IV.subsubsec6.8.5]{6.8.5}) 与 $\rho$ 的 pro-伴随 $\sigma$ 的复合).
(注意: 对 $y\in Y$, $\phi(y)$ 可与\emph{由 $\rho$ 定义的粘合层}
\eqhyperref[IV.eq9.6.4]{9.6.4} \emph{在 $y$ 处的纤维}等同.)
证明若 $E$ 等价于形如 $\Top(Z)$ 的拓扑斯,
则 $\phi$ 的值落在 $\Pro(\Ouv(X))$ 的本质像中.
\end{enumerate}
\end{enonce*}

\begin{enonce*}[remark]{问题}
逆命题是否成立?
\end{enonce*}

% END INCLUDED SGA4_English_current_buildable_source_20260906/source/04/drafts/SGA4_IV_LINES_7244_7443_English.texfrag


% BEGIN INCLUDED SGA4_English_current_buildable_source_20260906/source/04/drafts/SGA4_IV_LINES_7444_7700_English.texfrag
\subsection{粘合层}\etiquette[9.6]{IV.subsec9.6}

我们在 \hyperref[IV.theo9.5.4]{9.5.4} 已看到,
给定一个配有开子对象 $U$ 的拓扑斯 $E$,
实质上等于给定两个拓扑斯 $\mathcal{U}$ 和 $\mathcal{F}$ 及一个粘合函子
\begin{equation*}
\rho:\mathcal{U}\longrightarrow \mathcal{F}\tag{9.6.1}
\end{equation*}\etiquette[9.6.1]{IV.eq9.6.1}
\ie 一个有 pro-伴随
\begin{equation*}
\sigma:\Pro(\mathcal{F})\longrightarrow
\Pro(\mathcal{U}){\quoi}.\tag{9.6.2}
\end{equation*}\etiquette[9.6.2]{IV.eq9.6.2}
的函子.
由关于 pro-伴随函子的成串推论 (\hyperrefext[I][I.sec8]{8}),
只要知道函子 $\sigma$,
或等价地, 知道由 $\sigma$ 诱导的函子
\begin{equation*}
\sigma_{\circ}:\mathcal{F}\longrightarrow
\Pro(\mathcal{U}){\quoi}.\tag{9.6.3}
\end{equation*}\etiquette[9.6.3]{IV.eq9.6.3}
就可在相差唯一同构的意义下确定函子 $\rho$
(从而在相差拓扑斯等价的意义下确定粘合拓扑斯
$E=(\mathcal{F},\mathcal{U},\rho)$).
函子 $\sigma$ 在相差唯一同构的意义下,
由它与小的滤过投射极限交换且延拓函子 $\sigma_\circ$ 这两个性质确定.
另一方面, 给定函子 \eqhyperref[IV.eq9.6.3]{9.6.3} 同构于某个 pro-伴随的充要条件,
显然是对 $\mathcal{U}$ 的每个对象 $X$,
$\mathcal{F}$ 上的反变函子
$$
Y\longmapsto \Hom_{\Pro(\mathcal{U})}(\sigma_{\circ}(Y),X)
$$
可表\pageoriginale,
或等价地, 它是 $\mathcal{F}$ 的典范拓扑的一个层
(\sisi{\hyperref[IV.subsubsec1.1.2]{1.1.2} (iii)}{\hyperref[IV.theo1.2]{1.2} iii)}).
立即可知, 这也等于说反函子
\begin{equation*}
\sigma_{\circ}^{\circ}:\mathcal{F}^{\circ}\longrightarrow
\Pro(\mathcal{U})^{\circ}\subset
\check{\mathcal{U}}=\Hhom(\mathcal{U},(\Ens))\tag{9.6.4}
\end{equation*}\etiquette[9.6.4]{IV.eq9.6.4}
是 $\mathcal{F}$ 上的一个取值于 $\Pro(\mathcal{U})^{\circ}$ 的\emph{层}
(\hyperrefext[II][II.def6.1]{6.1}).
因此, \emph{给定一个粘合函子} \eqhyperref[IV.eq9.6.1]{9.6.1},
\emph{也等价于} (在相差同构的意义下)
\emph{给定 $\mathcal{F}$ 上的一个取值于} $\Pro(\mathcal{U})^{\circ}$ \emph{的层}.
与粘合函子 (或 (\hyperref[IV.subsec9.3]{9.3}) 中的情形 $(E,U,\mathcal{F})$)
如此相应的层,
称为与所考虑的粘合函子 $\rho$
(\resp 与拓扑斯 $E$ 的开子对象 $U$) 相应的\emph{粘合层}.

\begin{enonce*}[remark]{习题 9.6.5}\etiquette[9.6.5]{IV.exer9.6.5}
将 \hyperref[IV.theo9.5.4]{9.5.4} 精确化为两个 2-范畴之间的 2-等价陈述:
一端是二元组 $(E,U)$ 所成的 2-范畴,
其中 $E$ 是一个拓扑斯, $U$ 是后者的一个开子对象;
另一端是三元组 $(\mathcal{U},\mathcal{F},\rho)$ 所成的 2-范畴,
其中 $\mathcal{U}$ 和 $\mathcal{F}$ 是拓扑斯,
$\rho:\mathcal{U}\rightarrow\mathcal{F}$ 是一个粘合函子
(所讨论的 2-范畴的显式描述留给读者).
给出这个陈述的一个变式,
其中使用三元组 $(\mathcal{U},\mathcal{F},\phi)$ 所成的 2-范畴,
此时 $\phi$ 是 $\mathcal{F}$ 上取值于 $\Pro(\mathcal{U})^{\circ}$ 的一个层.
\end{enonce*}

\subsection{粘合拓扑斯的点}\etiquette[9.7]{IV.subsec9.7}

\subsubsection{}\etiquette[9.7.1]{IV.subsec9.7.1}
设 $C$ 是一个范畴,
$$
u:D'\longrightarrow C\quav v:D''\longrightarrow C
$$
是两个全忠实函子,
使 $u$ 和 $v$ 的本质像 $C'$ 和 $C''$ 满足
$$
\Ob C'\cap \Ob C''=\phi\quav \Ob C'\cup \Ob C''=\Ob C{\quoi}.
$$
再假设
$$
X'\in \Ob C'\quav X''\in \Ob C''\Rightarrow
\Hom(X'',X')=\phi{\quoi},
$$
或等价地,
$$
\Hom(v(A''), u(A'))=\phi\mbox{ 对 } A'\in \Ob D', A''\in \Ob
D''{\quoi}.
$$
于是显然, 由范畴 $C'$ 和 $C''$ 以及函子
$$
(X',X'')\longmapsto \Hom(X',X''):{C'}^{\circ}\times
C''\longrightarrow (\Ens){\quoi};
$$
可在相差典范\pageoriginale 同构的意义下重构范畴 $C$;
类似地, 由 $D'$, $D''$ 以及函子
$$
(A',A'')\longmapsto \Hom(u(A'), v(A'')):{D'}^{\circ}\times
D''\longrightarrow (\Ens){\quoi}.
$$
可在相差等价的意义下重构 $C$
(这个等价在相差唯一同构的意义下有定义).
这个观察尤其适用于下一个命题所述的情形,
并表明范畴 $\Ppoint(E)$
在相差等价的意义下由范畴 $\Ppoint(\mathcal{U})$ 和 $\Ppoint(\mathcal{F})$
以及某个从粘合函子 $\rho:\mathcal{U}\rightarrow\mathcal{F}$ 导出的函子
$$
\Ppoint(\mathcal{U})^{\circ}\times
\Ppoint(\mathcal{F})\longrightarrow (\Ens)
$$
确定.

\begin{enonce*}{命题 9.7.2}\etiquette[9.7.2]{IV.prop9.7.2}
设 $E$ 是一个拓扑斯, $U$ 是 $E$ 的一个开子对象,
$\mathcal{U}=E_{/U}$,
$\mathcal{F}$ 是 $E$ 中与 $U$ 互补的闭子拓扑斯.
考虑由包含态射 $j:\mathcal{U}\rightarrow E$
和 $i:\mathcal{F}\rightarrow E$ 诱导的函子
$p\mapsto j\circ p$ 和 $q\mapsto i\circ q$:
\begin{equation*}
u:\Ppoint(\mathcal{U})\longrightarrow \Ppoint(E),\quad
v:\Ppoint(\mathcal{F})\longrightarrow
\Ppoint(E){\quoi}.\tag{9.7.2.1}
\end{equation*}\etiquette[9.7.2.1]{IV.eq9.7.2.1}
则以下断言成立:
\begin{enumerate}
\renewcommand{\theenumi}{\alph{enumi}}
\renewcommand{\labelenumi}{\theenumi)}
\item 上述函子 $u$ 和 $v$ 全忠实,
且 $E$ 的每个点恰好属于函子 $u$ 和 $v$ 中一个的本质像.

\item 设 $p$ 和 $q$ 是 $E$ 的两个点,
使 $p$ (\resp $q$) 属于 $u$ (\resp $v$) 的本质像. 则
\begin{equation*}
\Hom(q,p)=\phi{\quoi}.\tag{9.7.2.2}
\end{equation*}\etiquette[9.7.2.2]{IV.eq9.7.2.2}

\item 设 $p$ (\resp $q$) 是 $\mathcal{U}$ (\resp $\mathcal{F}$) 的一个点.
则有一个关于 $(q,p)$ 具有函子性的同构:
\begin{equation*}
\Hom(u(p),v(q))\simeq \Hom(\Pro(\rho)(p),q)\simeq
\Hom(p,\sigma(q)){\quoi},\tag{9.7.2.3}
\end{equation*}\etiquette[9.7.2.3]{IV.eq9.7.2.3}
其中
$$
\sigma:\Pro(\mathcal{F})\longrightarrow \Pro(\mathcal{U})
$$
是粘合函子 $\rho:\mathcal{U}\rightarrow\mathcal{F}$ 的 pro-伴随函子
$\eqhyperref[IV.eq9.6.2]{9.6.2}$;
这里在写
$$
\rho:\mathcal{U}\longrightarrow \mathcal{F}{\quoi},
$$
时, 我们\pageoriginale 已在相差等价的意义下,
分别将范畴 $\Ppoint(\mathcal{U})$ 和 $\Ppoint(\mathcal{F})$
与 $\Pro(\mathcal{U})$ 和 $\Pro(\mathcal{F})$ 的全子范畴等同
(\hyperref[IV.subsubsec6.8.5]{6.8.5}).
\end{enumerate}
\end{enonce*}

\noindent \emph{证明}
\begin{enumerate}
\renewcommand{\theenumi}{\alph{enumi}}
\renewcommand{\labelenumi}{\theenumi)}
\item $u$ 和 $v$ 全忠实是已知的
(\hyperref[IV.prop9.1.4]{9.1.4}).
关于 $u$ 和 $v$ 的本质像的断言,
由点 $P\rightarrow E$ 在相差同构的意义下
经过 $\mathcal{U}$ \resp $\mathcal{F}$ 分解的判据
\hyperref[IV.subsubsec9.2.5]{9.2.5} 和 \hyperref[IV.prop9.4.2]{9.4.2} 得出;
这里用到点拓扑斯 $P$ 仅有两个开子对象,
即 $P$ 的始对象和终对象.

\item 根据 a), 这个断言就是说:
若点 $p:P\rightarrow E$ 经过 $\mathcal{U}$ 分解,
则每个满足 $\Hom(p',p)\neq\phi$ 的点 $p':P\rightarrow E$ 也如此.
现在若 $\alpha:p'\rightarrow p$,
则它诱导 $p^{\ast}(U)\rightarrow{p'}^{\ast}(U)$,
而 $p^{\ast}(U)=e_p$ 便蕴含 ${p'}^{\ast}(U)=e_p$, \cqfd

\item 将下文引理 \hyperref[IV.lem9.7.2.4]{9.7.2.4}
应用于包含 $j:\mathcal{U}\rightarrow E$ 和 $i:\mathcal{F}\rightarrow E$,
得到一个在相差典范同构的意义下交换的函子图:
\[
\xymatrix@C=2cm@R=.55cm@M=.18cm{ \Point(\Uu) \ar@{^{(}->}[d]
\ar[r]^{u}& \Point(E) \ar@{^{(}->}[d]&
\Point(\Ff) \ar@{^{(}->}[d] \ar[l]_{v}\\
\Pro(\Uu) \ar[r]^{\Pro(j_{\ast})}  & \Pro(E)   & \Pro(\Ff)
\ar[l]_{\Pro(i_{\ast})} , }
\]
其中竖箭头表示典范全忠实函子
(\hyperref[IV.subsubsec6.8.5]{6.8.5}).
此外, 由于 $i_{\ast}$ 有左伴随 $i^{\ast}$,
$\Pro(i_{\ast})$ 有左伴随 $\Pro(i^{\ast})$
\hyperrefext[I][I.sec8.11.5]{8.11.5},
最后得到函子同构
\begin{align*}
\Hom(u(p),v(q)) \simeq & \Hom_{\Pro(E)}(\Pro(j_{\ast})(p),
\Pro(i_{\ast})(q))\\
& \simeq \Hom_{\Pro(\mathcal{F})}(\Pro(i^{\ast})(\Pro(j_{\ast})(p)),q)\\
& \simeq \Hom_{\Pro(\mathcal{F})}(\Pro(i^{\ast}j_{\ast})(p),q),
\end{align*}
这正是 \hyperref[IV.eq9.7.2.3]{(9.7.2.3)} 的第一个公式;
第二个公式直接由 $\sigma$ 作为 $\rho$ 的 pro-伴随的定义得出.
这就证明了 \hyperref[IV.prop9.7.2]{9.7.2}, 只待证明以下引理.
\end{enumerate}

\begin{enonce*}{引理 9.7.2.4}\etiquette[9.7.2.4]{IV.lem9.7.2.4}
设 $f:F\rightarrow E$ 是一个拓扑斯态射. 则函子图
\[\vcenter{
\xymatrix@C=2cm@R=.5cm{
  \Point(F)\ar[r]^{p \mapsto f \circ p} \ar[d] & \Point(E)\ar[d]\\
\Pro(F) \ar[r]^{\Pro(f_{\ast})}     & \Pro(E) }}\leqno{(9.7.2.4.1)}
\]\etiquette[9.7.2.4.1]{IV.eq9.7.2.4.1}
在相差典范同构的意义下\pageoriginale 交换
(竖箭头表示全忠实函子 $\hyperref[IV.subsubsec6.8.5]{6.8.5}$).
\end{enonce*}

由定义直接可验证, 验证留给读者.

\begin{enonce*}{推论 9.7.3}\etiquette[9.7.3]{IV.cor9.7.3}
设 $E$ 是一个有足够多点的拓扑斯.
则 $E$ 的每个闭子拓扑斯 $\mathcal{F}$ 也有足够多的点.
更确切地, 若 $(p_i)_{i\in I}$ 是 $E$ 的一个保守点族,
则由其中来自 $\mathcal{F}$ 的点的那些 $p_i$ 所成的子族 $(p_i)_{i\in J}$,
所定义的 $\mathcal{F}$ 的点族保守.
(比较 $\hyperref[IV.cor6.7.3]{6.7.3}$, $\hyperref[IV.rem6.7.4]{6.7.4}$.)
\end{enonce*}

事实上, 设 $f:X\rightarrow Y$ 是 $\mathcal{F}$ 的一个箭头,
在所考虑的 $\mathcal{F}$ 的点所对应的纤维函子下都映为同构.
由于 $i^{\ast}i_{\ast}\simeq\id_{\mathcal{F}}$,
可知 $i_{\ast}(f):i_{\ast}(X)\rightarrow i_{\ast}(Y)$
在与 $i\in J$ 的 $p_i$ 相应的纤维函子下都映为同构;
对于 $i\in I-J$ 的 $p_i$, \ie 来自 $U$ 的那些点, 也如此,
因为 $j^{\ast}i_{\ast}$ 是以终对象为值的常值函子.
因此 $i_{\ast}(f)$ 是同构,
从而 $f\simeq i^{\ast}i_{\ast}(f)$ 也是同构, \cqfd

\begin{enonce*}[remark]{注 9.7.4}\etiquette[9.7.4]{IV.rem9.7.4}
若 $E$ 是一个有足够多点的拓扑斯,
显然只要知道 $\Ppoint(E)$ 中作为 $\Ppoint(E_{/U})$ 的本质像的全子范畴,
就确定了 $E$ 的开子对象 $U$;
事实上, 若 $C$ 是定义 $E$ 的一个保守点族的 $\Ppoint(E)$ 的全子范畴,
则甚至只须知道由 $C$ 中属于 $\Ppoint(U)$ 的本质像的元素
所成的子范畴 $C_{/U}$.
由此及 \hyperref[IV.prop9.7.2]{9.7.2} a) 可知,
只要知道 $\Ppoint(\mathcal{F})$ 在 $\Ppoint(E)$ 中的本质像,
甚至仅知道它与 $C$ 的交,
就确定了 $E$ 的一个闭子拓扑斯 $\mathcal{F}$.
\end{enonce*}

\begin{enonce*}[remark]{习题 9.7.5}\etiquette[9.7.5]{IV.exer9.7.5}

% END INCLUDED SGA4_English_current_buildable_source_20260906/source/04/drafts/SGA4_IV_LINES_7444_7700_English.texfrag


% BEGIN INCLUDED SGA4_English_current_buildable_source_20260906/source/04/drafts/SGA4_IV_LINES_7701_7780_English.texfrag
设 $E$ 是一个拓扑斯.
对 $E$ 的每个子拓扑斯 $F$,
设 $P(F)$ 是 $\Point(E)$ 中由经过 \sisi{$E$}{$F$} 分解的 $E$ 的点的同构类组成的子集.
于是 $P(F)$ 同胚于 $\Point(F)$ (\hyperref[IV.exer9.1.8]{9.1.8} c)).
证明若 $F$ 是 $E$ 的一个开
(\resp 闭, \resp 局部闭 (\hyperref[IV.exer9.4.9]{9.4.9})) 子拓扑斯,
则 $P(F)$ 是 \sisi{$E$}{$\Point(E)$} 的开\pageoriginale
(\resp 闭, \resp 局部闭) 子集.
证明当 $E$ 有足够多的点时,
 $E$ 的每个局部闭子拓扑斯也如此,
且映射 $F\mapsto P(F)$
诱导 $E$ 的开 (\resp 闭, \resp 局部闭) 子拓扑斯所成的集合,
与 $\Point(E)$ 的开 (\resp 闭, \resp 局部闭) 子集所成的集合之间的一个双射.
将上述讨论推广到把 $\Point(E)$ 换成 $\Point(E)$ 的任意一个与 $E$ 的保守点族对应的子空间 $X$ 的情形.
\end{enonce*}

\subsection{关于拓扑空间所对应的某些拓扑斯的补充}\etiquette[9.8]{IV.subsec9.8}

本节以习题的形式给出一些补充.
建议读者浏览这些习题, 以熟悉在各种经典情形中采用拓扑斯观点的方法.

\begin{enonce*}[remark]{习题 9.8.1}\etiquette[9.8.1]{IV.exer9.8.1}
(\emph{拓扑空间上的层的下降.})

设 $f:X\rightarrow Y$ 是 $𝒰$ 中的拓扑空间之间的一个连续映射,
给出一个函子
$$
f^{\ast}:\Top(Y)\longrightarrow\Top(X){\quoi}.
$$
\begin{enumerate}
\renewcommand{\theenumi}{\alph{enumi}}
\renewcommand{\labelenumi}{\theenumi)}
\item 以下条件等价:
\begin{enumerate}
\renewcommand{\theenumii}{\roman{enumii}}
\renewcommand{\labelenumii}{(\theenumii)}
\item $f^{\ast}$ 忠实.

\item $f^{\ast}$ 保守.

\item $f(X)$ 是 $Y$ 的一个很稠密子集 (EGA IV 10.1.3).
\end{enumerate}
$f$ 或 $f_{\sob}$ 满射就足够了.
给出一个 $f^{\ast}$ 忠实而 $f=f_{\sob}$ 不满射的例子
(取 $X$ 离散, $f$ 单射, $f(X)$ 为 $Y$ 的闭点集,
$Y$ 为一个非离散的 Noether 拓扑空间).
给出一个 $f_{\sob}$ 满射而 $f$ 不满射的例子,
再给出一个 $f$ 满射而 $f_{\sob}$ 不满射的例子.

\item 假设 $f(X)$ 在 $Y$ 中很稠密,
且其拓扑是从 $X$ 诱导的商拓扑.
证明 $(\Esp)$ 的箭头 $f$
对于变动空间上的集合层所成的纤维化范畴是一个下降态射,
\ie 函子 $f^{\ast}$
诱导\pageoriginale 一个从 $\Top(Y)$
到配备相对于 $f:X\rightarrow Y$ 的下降数据的
$\Top(X)$ 的对象所成范畴的\emph{全忠实}函子.
归结到 $f$ 满射的情形,
并仿照 VIII \hyperref[VIII.prop9.1]{9.1} 中针对概形平展拓扑给出的论证.
给出一个逆命题.

\item 假设 $f(X)$ 在 $Y$ 中很稠密,
其拓扑是从 $X$ 诱导的商拓扑, 且 $f$ 的纤维连通.
证明函子 $f^{\ast}$ 全忠实.
(仿照 \SGA 1 IX 3.4 的论证.)
至少在 $Y$ 的点均为闭点的情形中, 给出一个逆命题.

\item 假设 $Y$ 是开子集 $Y_i$ 的并,
且 $X$ 在其上分别有截面.
证明此时 $f$ 对于变动拓扑空间上的集合层所成的纤维化范畴,
是一个\emph{有效下降态射},
\ie b) 中考虑的函子事实上是范畴等价.
\end{enumerate}
\end{enonce*}

% END INCLUDED SGA4_English_current_buildable_source_20260906/source/04/drafts/SGA4_IV_LINES_7701_7780_English.texfrag


% BEGIN INCLUDED SGA4_English_current_buildable_source_20260906/source/04/drafts/SGA4_IV_LINES_7781_8025_English.texfrag
\begin{enonce*}[remark]{习题 9.8.2}\etiquette[9.8.2]{IV.exer9.8.2}
(\emph{拓扑空间的商拓扑斯. 拓扑展拓扑斯（étendue）.})

\begin{enumerate}
\renewcommand{\theenumi}{\alph{enumi}}
\renewcommand{\labelenumi}{\theenumi)}
\item 设
\hypertarget{sga4.exp.IV.equation.9.8.2.1}{}\[
\xymatrix@C=1.5cm@R=1.5cm@M=.2cm{
X_2 \ar@<6pt>[r] \ar@<-1.5pt>[r] \ar@<-8pt>[r]& X_1 \ar@<2pt>[r]^{p_1} \ar@<-4pt>[r]_{p_2}& X_0\\
}\leqno{(9.8.2.1)}
\]\etiquette[9.8.2.1]{IV.eq9.8.2.1}
是在拓扑空间 $\in𝒰$ 所成的范畴 $(\Esp)$ 中,
一个截断到次数 $2$ 的半单纯对象;
逆像函子于是给出一个层范畴的图
\[
\xymatrix@M=.2cm@C=1.25cm{
  \Top(X_0) \ar@<2pt>[r]^{p_1^{\ast}} \ar@<-2pt>[r]_{p_2^{\ast}}&
  \Top(X_1) \ar@<-5pt>[r] \ar[r] \ar@<5pt>[r]& \Top(X_2).
}
\]
(更确切地, 有一个在标准单形 $\Delta_n$, $0\leq n\leq2$ 的范畴上余纤维化的范畴.)
证明这个图的 $\varprojlim$ 范畴
(\ie $X_0$ 上\emph{配备}相对于图
\sisi{(9.8.6.1)}{\hyperref[IV.eq9.8.2.1]{9.8.2.1}} 的下降数据的集合层范畴,
\ie 配备满足\ldots 的同构 $p^{\ast}_1(F)\simeq p^{\ast}_2(F)$)
是一个 $𝒰$-拓扑斯 $T$.
(关于拓扑斯在 $\varprojlim$ 运算下稳定性的更一般陈述,
\cf \sisi{VI)}{\hyperrefext[VI][VI.sec8]{8}}.)
定义一个拓扑斯态射 $q:\Top(X_0)\rightarrow T$,
并证明对每个 $𝒰$-拓扑斯\pageoriginale $E$,
$\Hhomtop(T,E)$ 经 $q^{\ast}$ 等价于以下范畴图的 $\varprojlim$ 范畴,
这个图由 \hyperref[IV.eq9.8.2.1]{(9.8.2.1)} 得出:
\[
\xymatrix@M=.2cm{ \Hhomtop(\Top(X_0),E)\ar@<2pt>[r] \ar@<-4pt>[r]&
\Hhomtop(\Top(X_1),E)\ar[r]\ar@<6pt>[r]\ar@<-6pt>[r]&
\Hhomtop(\Top(X_2),E). }
\]
\item 特别地, 设 $R$ 是 $(\Esp)$ 的一个对象 $X$ 中的等价关系,
\ie (\sisi{\hyperrefext[I][I.sec10]{10}}{\hyperrefext[I][I.sec10.9]{10.9}})
$X\times X$ 的一个子对象 $R$
(注意包含 $R\hookrightarrow X\times X$ 连续,
但 $R$ 的拓扑未必由 $X\times X$ 的拓扑诱导),
满足对 $(\Esp)$ 的每个对象 $Y$,
$\Hom(Y,R)$ 是 $\Hom(Y,X)$ 中一个等价关系的图.
由此得到形如 \hyperref[IV.eq9.8.2.1]{(9.8.2.1)} 的图,
其中 $X_0=X$, $X_1=R$,
$X_2=\fprod{(R,p_2)}{R,p_1}{X}$,
$p_1$, $p_2$ 是从 $R$ 到 $X$ 的两个投射;
从而得到一个拓扑斯,
我们将它记为 $\Top(X)/R$,
甚至滥用记号写为 $\Top(X/R)$ 或 $X/R$,
它在 a) 所明确的意义下起到 $X$ 对 $R$ 的商的作用.
假设 $R$ 的拓扑由 $X\times X$ 的拓扑诱导,
并用 $X/R$ 表示通常的商拓扑空间,
假设 $X$ 在 $X/R$ 上局部有截面.
证明此时 $T$ 等价于 $\Top(X/R)$.
(利用 \hyperref[IV.exer9.8.1]{9.8.1} d).)

\item 设 $H\rightarrow X$ 是拓扑群的一个单射态射,
\ie $(\Esp)$ 中群对象的一个单态射;
设 $R$ 是它在 $X$ 中定义的等价关系,
\ie $R=H\times X$, 其中 $p_1=\pr_2$,
$p_2$ 由 $H$ 在 $X$ 上的左平移作用定义.
证明 $R$ 的拓扑由 $X\times X$ 的拓扑诱导,
当且仅当 $H$ 的拓扑由 $X$ 的拓扑诱导.
在以下两种情形中给出不满足这个条件,
并且 $X/H$ 上的通常商拓扑为平庸拓扑的例子:
a) $H=\mathbb{Z}\times\mathbb{Z}$, $X=\mathbb{R}$;
b) $H=\mathbb{R}$, $X=\mathbb{T}\times\mathbb{T}$,
其中 $\mathbb{T}=\mathbb{R}/\mathbb{Z}$ (“环面上的测地线”).
证明这样得到的两个拓扑斯等价,
且既不等价于 $\Top(X/H)$
(它是终拓扑斯 (\hyperref[IV.subsec2.2]{2.2})),
也不等价于任何形如 $\Top(Y)$ 的拓扑斯.

\item 设 $G$ 是作用在拓扑空间 $X$ 上的一个拓扑群,
由此按熟悉的方式得到一个半单纯对象
\[
\xymatrix@M=.2cm{ ... G\times G \times X
\ar@<6pt>[r]\ar@<-1.5pt>[r]\ar@<-8pt>[r] &
  G\times X \ar@<3pt>[r] \ar@<-4pt>[r] & X.
}
\]
证明\pageoriginale 由 a) 从它得出的拓扑斯 $T$,
可与如下空间 $X'$ 所成的范畴等同:
它们在 $X$ 上, 以 $G$ 为算子群,
并且 $X'\rightarrow X$ 是一个平展映射 (与 $G$ 的作用相容).
特别地, 当 $G$ 离散时,
重新得到 \hyperref[IV.subsec2.3]{2.3} 中的拓扑斯 $\Top(X,G)$.

上述讨论说明, 可以将这个拓扑斯 $T$
记为 $\Top(X)/G$, 甚至 $\Top(X/G)$, 或简记为 $X/G$.
但应记住: 当离散群 $G$ (为明确起见) \emph{固有地}作用在 $X$ 上,
从而通常的商空间 $X/G$ 有相当好的性质时,
由 $T$ 的泛刻画 a) 得出的自然拓扑斯态射 $T\rightarrow\Top(X/G)$,
只有在 $G$ 自由作用,
\ie 没有不动点, \ie $G\times X\rightarrow X\times X$ 是单态射时,
才是拓扑斯等价.
因此, 若一个“预等价关系”
(或 (\SGA 3 V 1) 意义下的“群胚”) 不是等价关系,
则在拓扑斯意义下取商 (或“取精细商”),
一般不通过对应 $Y\mapsto\Top(Y)$ 与通常的拓扑商相对应.

\item 设 $T$ 是一个 $𝒰$-拓扑斯. 证明以下条件等价:
\begin{itemize}
\item[(i)] 存在一个覆盖 $T$ 的终对象的 $T$ 的对象族 $(S_i)_{i\in I}$,
使对每个 $i\in I$,
诱导拓扑斯 $T_{/S_i}$ 等价于形如 $\Top(X)$ 的拓扑斯,
其中 $X\in(\Esp)$.

\item[${\rm (i')}$] 存在一个覆盖终对象的 $S\in\Ob T$,
使诱导拓扑斯 $T_{/S}$ 等价于形如 $\Top(X)$ 的拓扑斯.

\item[(ii)] 存在一个拓扑空间 $X$ 及 $X$ 中的一个\emph{平展}预等价关系 $R$
(在范畴 $(\Esp)$ 的意义下)
(\ie $p_1:R\rightarrow X$ 是平展映射),
使采用 b) 的记号时 $T$ 等价于 $\Top(X/R)$.
\end{itemize}

此时\pageoriginale 称拓扑斯 $T$ \emph{局部为拓扑空间},
或称为\emph{拓扑展拓扑斯}
(无混淆可能时, 简称\emph{展拓扑斯}).

\item 证明 a) 中构造的拓扑斯 $T$ 有足够多的点;
更确切地, 它有一个由 (小) 集合 $X_0$ 参数化的保守点族.
特别地, 每个展拓扑斯都有一个小保守点族, 从而有足够多的点.
当展拓扑斯 $T$ 如 e) ii) 那样表示为 $\Top(X/R)$ 时,
证明 $T$ 的每个点都同构于 $\Top(X)$ 的一个点的像,
从而由 $X_{\sob}$ 的一个普通点定义
(当 $X$ 清醒时, 由 $X$ 自身的一个普通点定义).

\item 证明 \hyperref[IV.subsec2.3]{2.3} 中的拓扑斯 $\Top(X/G)$ 是展拓扑斯.
(利用 d) 中的描述或 \hyperref[IV.exer7.1.10]{7.1.10} c).)
证明对每个 $x\in X$,
由 $x$ 定义的 $\Top(X,G)=\Top(X)/G$ 的点的自同态幺半群,
典范同构于 $x$ 的稳定子群 $G_x$.
特别地, 展拓扑斯的点可以有非平凡自同构群.
对于由空间 $X$ 中的一个平展预等价关系定义的展拓扑斯 $T$,
确定 $\Ppoint(T)$,
并证明 $T$ 的点的每个自同态都是自同构.

\item 设 $T$ 是一个拓扑斯. 证明以下条件等价:
\begin{itemize}
\item[i)] 对 $T$ 的每个点 $p$,
$p$ 的每个自同态 (分别 $p$ 的每个自同构) 都是恒等.

\item[ii)] 对每个有足够多点的拓扑斯 $S$
以及每个拓扑斯态射 $f:S\rightarrow T$,
$f$ 的每个自同态 (分别 $f$ 的每个自同构) 都是恒等.
\end{itemize}
同样, 证明以下条件等价:
\begin{itemize}
\item[${\rm i')}$] 对 $T$ 的两个点 $p$, $q$,
从 $p$ 到 $q$ 至多有一个态射.

\item[${\rm ii')}$] 对从有足够多\pageoriginale 点的拓扑斯 $S$ 到 $T$
的两个态射 $f$, $g:S\rightrightarrows T$,
从 $f$ 到 $g$ 至多有一个态射.
\end{itemize}

当后面这些条件成立时,
称拓扑斯 $T$ \emph{逐点刚性}, 或简称\emph{刚性}.
证明若拓扑空间 $X$ 配有一个平展预等价关系 $R$,
则商展拓扑斯 $T=\Top(X)/R$ 刚性,
当且仅当 $R$ 是一个等价关系.

\item 设 $(X,G)$ 是一个带离散算子群的拓扑空间,
$H$ 是 $G$ 的一个正规子群,
使 $H$ 在 $X$ 上的诱导作用固有且自由;
设 $X'=X/H$, $G'=G/H$,
于是 $G'$ 按显然的方式作用在 $X'$ 上.
证明由带算子空间的典范态射 $(X,G)\rightarrow(X',G')$
(\hyperref[IV.subsubsec4.1.2]{4.1.2}) 诱导的拓扑斯态射
$$
\Top(X,G)\longrightarrow \Top(X',G')
$$
是一个拓扑斯等价.

证明 $\Top(X/G)$ 的开子拓扑斯所成的集合,
可与 $X$ 中在 $G$ 的作用下稳定的开子集所成的集合等同,
或等价地, 与商拓扑空间 $X/G$ 的开子集集合等同.
由此得出, $\Top(X/G)$ 连通,
当且仅当 $X$ 中每个既开又闭且在 $G$ 的作用下稳定的子集,
不是空集就是整个 $X$.
当 $X$ 的连通分支都是开集时,
$\Top(X/G)$ 连通, 当且仅当 $G$ 在 $X$ 的连通分支集合 $\pi_0(X)$ 上传递作用;
更一般地, $\pi_0(\Top(X/G))$
(\hyperref[IV.exer8.7]{8.7} g))
是一个本质常值的 pro-集合, 同构于商集 $\pi_0(X)/G$.

证明当 $X$ 局部连通且局部单连通,
而 $T=\Top(X,G)$ 连通,
\ie $G$ 在 $\pi_0(X)$ 上传递作用时,
在通过带算子空间 $(X',G')$ 表示 $T$ 的所有方式中
(上文给出了一些这种表示),
有一种使 $X'$ 连通且单连通,
并且这种表示在相差不唯一同构的意义下唯一.
证明选取这样的 $(X',G')$,
等于选取 $T$ 的终对象的一个“万有覆盖”,
并且 $G'$ 同构于关于此万有覆盖的群 $\pi_1(T)$
(\cf \hyperref[IV.exer2.7.5]{2.7.5}).

\noindent (提示:\pageoriginale
证明给定 $T$ 与某个 $\Top(X,G)$ 的一个等价,
等于给定 $T$ 中的一个 $G_T$-挠子 $P$;
并且若 $(X,G)$ 与 $P$ 相应,
则有典范等价 $\Top(X)\simeq T_{/P}$.)

由此给出 c) 的最后一个断言的显然证明.
将以下拓扑斯 $T$ 刻画为连通、局部连通且局部单连通的展拓扑斯:
它们等价于某个 $\Top(X,G)$,
其中 $G$ 是作用在局部连通且局部单连通拓扑空间上,
并传递地置换其连通分支的离散群;
这个刻画还要求 $T$ 的终对象 $e_T$ 的万有覆盖 $P$
满足诱导拓扑斯 $T_{/P}$ 等价于形如 $\Top(X)$ 的拓扑斯
(或如我们将滥用语言简说的,
$T_{/P}$ “是”一个拓扑空间).

给出一个局部同构于 $\mathbb{R}$ 的展拓扑斯的例子,
使它连通且单连通, 却不是拓扑空间.
(取有双原点的直线 $\mathbb{R}$ 的万有覆盖.)

\item 若环化拓扑斯 $(T,{\Oo}_T)$ (\cf \hyperref[IV.sec11]{\S 11})
满足以下条件, 就称它是\emph{可微展拓扑斯}
(分别\emph{实解析展拓扑斯}, 分别\emph{复解析展拓扑斯}):
存在覆盖终对象的 $T$ 的对象 $S_i$,
使诱导环化拓扑斯 $(T_{/S_i},{\Oo}_T/S_i)$
等价于由可微流形及其取实值的 $C^{\infty}$ 函数层定义的环化拓扑斯
(分别等价于由实解析空间定义的环化拓扑斯,
分别等价于由复解析空间定义的环化拓扑斯).
给出这些环化拓扑斯的 e) ii) 型的构造性描述
(分别取 $X$ 为可微流形、实解析空间或复解析空间).
回到 c) 中考虑的例子,
给出这类环化拓扑斯的例子,
使它们不等价于拓扑空间所对应的环化拓扑斯.
\end{enumerate}
\end{enonce*}

% END INCLUDED SGA4_English_current_buildable_source_20260906/source/04/drafts/SGA4_IV_LINES_7781_8025_English.texfrag


% BEGIN INCLUDED SGA4_English_current_buildable_source_20260906/source/04/drafts/SGA4_IV_LINES_8026_8245_English.texfrag
\begin{enonce*}[remark]{习题 9.8.3}\etiquette[9.8.3]{IV.exer9.8.3}
(\emph{与局部等价关系和叶状结构相应的拓扑斯\pageoriginale})
\footnote[1]{这道习题是在仓促之间写成, 且未经充分核查,
因此提出时保留一切意见.
N. \textsc{Saavedra} 先生核查了 a) 至 e) 部分.
请读者将任何意见告知我们.}.

设 $X$ 是一个拓扑空间.
\begin{enumerate}
\renewcommand{\theenumi}{\alph{enumi}}
\renewcommand{\labelenumi}{\theenumi)}
\item 对 $X$ 的每个开子集 $U$,
设 $\Quot(U)$ 是 $U$ 上的等价关系所成的集合;
对开子集 $V\subset U$, 考虑自然映射 $\Quot(U)\rightarrow\Quot(V)$.
这给出 $X$ 上的一个预层 $\Quot_X$.
设 $Q_X$ 或简记 $Q$ 为相伴层,
$r$ 是 $Q$ 的一个截面 (“\emph{$X$ 上的局部等价关系}”).
由 $X$ 的开子集 $U$ 上的等价关系集合 $\Quot(U)$ 上的序关系,
在预层 $\Quot$ 以及相伴层 $Q$ 上导出一个序结构,
\ie 一个取值于有序集范畴的预层 \resp 层结构.

\item 对 $X$ 上的每个等价关系 $R$,
考虑它定义的 $Q$ 的截面 $\loc(R)$.
设 $E(r)$ 是 $X$ 上使 $\loc(R)$ 比 $r$ 更粗的等价关系所成的集合,
$\glob(r)$ 是 $E(r)$ 中等价关系的最小上界
(它的图是所有 $R\in E(r)$ 的图的交).
设 $(U_x)_{x\in X}$ 是点 $x\in X$ 的一个开邻域族;
对每个 $x\in X$, 设 $R_x$ 是 $U_x$ 上的一个等价关系,
其在 $x$ 处的芽为 $r_x$.
对每个开子集族 $𝒱$, 其元素为 $V_x$ $(x\in X,x\in V_x\subset U_x)$,
设 $R_{𝒱}$ 为关系族 $R_x|V_x$ 生成的 $X$ 上的等价关系.
证明若 $𝒱'\leq𝒱$ (在显然的意义下),
则 $R_{𝒱'}\geq R_{𝒱}$,
且 $\glob(r)$ 是等价关系的递增滤过族 $R_{𝒱}$ 的最小上界.
给出一个 $\glob(r)$ 不属于 $E(r)$ 的例子,
\ie 存在一个 $a\in X$,
使对 $a$ 的每个\pageoriginale 开邻域 $W\subset U_a$,
存在 $𝒱=(V_x)$ 以及 $W$ 的两个点 $y,z$,
它们在 $R_a$ 下等价, 却不在 $R_{𝒱}$ 下等价.
证明对每个 $R\in\Quot(X)$, 有 $\glob\loc(R)\geq R$.

\item 若 $\glob(r)\in E(r)$,
\ie $\loc(\glob(r))\leq r$
(\resp 若对 $X$ 的每个开子集 $U$, $r$ 在 $U$ 上的限制均预凝聚),
就称局部等价关系 $r$ \emph{预凝聚（precoherent） (\resp 凝聚（coherent）)}.
若它凝聚, 并且还满足 $r=\loc\glob(r)$,
就称它整体凝聚.
若 $r=\loc(R)$ 凝聚, 就称 $X$ 上的等价关系 $R$ \emph{局部凝聚};
若还满足 $R=\glob(r)$,
\ie $R=\glob(\loc(R))$, 就称它\emph{凝聚}.
证明 $X$ 上整体凝聚的局部等价关系 $r$
与 $X$ 上凝聚的等价关系 $R$ 通过 $r\mapsto\glob(r)$
和 $R\mapsto\loc(R)$ 双射对应.

\item 证明 $r$ 凝聚的一个充分条件是:
对每个 $a\in X$ 以及 $a$ 的每个开邻域 $W'\subset U_a$,
存在 $a$ 的一个开邻域 $W\subset W'$,
使 $R_a|W$ 的每个等价类包含在 $R_a|W'$ 的某个等价类的一个连通分支中
(因此, 只须 $a$ 有一个开邻域基本系 $W_a\subset U_a$,
使诱导等价关系 $R_a|W_a$ 的纤维连通).
(提示: 归结到 $r$ 由某个 $R\in\Quot(X)$ 定义的情形,
并证明其预凝聚;
此时注意, 等价关系 $R_{𝒱}$ 的纤维是 $R$ 的纤维的相对开子集,
从而也是相对闭子集.)
证明一个整体等价关系 $R$ 凝聚的一个充分条件是:
它局部凝聚, 且其纤维连通;
证明若纤维闭, 则这个条件也是必要的.
给出一个纤维连通且局部连通, 却不局部凝聚的等价关系的例子 (?).

\item 设\pageoriginale $f:X'\rightarrow X$ 是一个连续映射.
定义一个自然预层同态
$f_{\ast}(\Quot_{X'})\leftarrow\Quot_X$,
使它诱导一个层同态 $f_{\ast}(Q_{X'})\leftarrow Q_X$,
从而诱导一个同态 $f^{\ast}(Q_X)\rightarrow Q_{X'}$.
若 $r$ 是 $X$ 上的一个局部等价关系,
并且 $r$ 的逆像是 $X'$ 上的粗局部等价关系,
就称 $f$ 是 $X$ 上局部等价关系 $r$ 的一个“\emph{纤维映射}”.
对每个拓扑空间 $X'$,
设 $\Homfib_r(X',X)$ 是从 $X'$ 到 $X$ 的连续映射中,
相对于 $r$ 为纤维映射的映射所成的集合.
证明当 $X'$ 变化时, 这定义了关于 $X'$ 的一个反变函子.

\item 假设 $r$ 凝聚.
证明上述函子可由 $X$ 上的一个拓扑空间 $X^r$ 代表.
证明万有纤维映射 $\Phi:X^r\rightarrow X$ 是双射,
并且每个 $x\in X^r$ 都有一个开邻域 $U$,
使由 $\Phi$ 诱导的映射 $U\rightarrow X$ 是 $U$ 到其像的同胚.
证明 $X^r$ 的连通分支都是开集,
并且在双射 $\Phi$ 下对应于等价关系 $R=\glob(r)$ 的纤维.
给出一个例子, 使 $\Phi$ 在 $X^r$ 的连通分支上的限制,
不诱导这些空间到它们在 $X$ 中的像的同胚.

\item 若局部等价关系 $r$ 是 $Q$ 的子层 $\Qouv_X$ 的一个截面,
就称它\emph{开};
其中这个子层来自 $\Quot_X$ 的子预层 $\Quotouv_X$,
后者在 $X$ 的每个开子集 $U$ 上的值是 $U$ 上的\emph{开}等价关系所成的集合.
若局部等价关系 $r$ 凝聚且开,
或等价地, 它凝聚, 且对 $X$ 的每个开子集 $U$,
$\glob(r|U)$ 是 $U$ 上的开等价关系,
就称它\emph{严格开}.
若 $X$ 上的等价关系 $R$ 凝聚,
且 $\loc(R)$ 是一个严格开的局部等价关系,
就称它\emph{严格开}.
证明开关系 $r$ 严格开的一个充分条件,
是它满足\pageoriginale d) 中考虑的条件;
若 $r$ 局部由某个具有闭纤维的 $R$ 定义,
则这个条件也是必要的.

\item 设 $F$ 是 $X$ 上的一个集合层.
对 $X$ 的每个开子集 $U$,
设 $\Quot(U,F)$ 是如下二元组所成的集合:
由 $U$ 上的一个等价关系 $R_U$
以及 $F|U$ 上的一个等价关系 $R_{F|U}$ 组成
(将 $F$ 视为 $X$ 上的一个平展空间),
使结构映射 $p:F|U\rightarrow U$
与等价关系 $R_U$ 和 $R_{F|U}$ 相容,
且相应拓扑空间图
\[
\xymatrix@C=2cm@R=1cm{
F|U \ar[r]\ar[d] \ar[d]_{p}& (F|U)/R_{F|U} \ar[d]_{q}\\
U   \ar[r]&  U/R_{U} }
\]
是 Cartesian 的, 而 $q$ 是平展映射
(从而 $F|U$ 可与 $U/R_U$ 上的层 $(F|U)/R_{F|U}$ 的逆像等同).
证明当 $U$ 变化时, $\Quot(U,F)$ 定义 $X$ 上的一个预层,
将其相伴层记为 $Q(F/X)$.
定义一个层同态 $Q(F/X)\rightarrow Q_X$.
对给定的 $Q_X$ 的截面 $r$,
层 $F$ 上的一个 $r$-\emph{结构},
是 $Q(F/X)$ 的一个位于 $Q_X$ 的给定截面 $r$ 上方的截面.
定义 $X$ 上的 $r$-层所成的范畴 $\Top(X/r)$,
并定义一个保守且忠实的“遗忘 $r$-结构”的函子
$\Top(X/r)\rightarrow\Top(X)$.
证明 $\Top(X/r)$ 中有限投射极限和有限归纳极限可表,
且上述函子与这些极限交换.
由此得出, $\Top(X/r)$ 中的有限和不交且普遍,
等价关系普遍有效.
证明 $\Top(X/r)$ 有一个小生成族.
给出一个 $r$ 凝聚而 $\Top(X/r)$ 没有无限余积的例子 (?),
从而它不是拓扑斯.

\item 设 $f:X\rightarrow Y$ 是一个与 $R=\glob(r)$ 相容的连续映射.
定义一个函子
$$
f^{\ast}_r:\Top(Y)\longrightarrow\Top(X/r)
$$
使它\pageoriginale 与包含函子 $\Top(X/r)\rightarrow\Top(X)$ 的复合为 $f^{\ast}$.
证明若 $r$ 整体凝聚且严格开
(\ie $R$ 严格开且 $r=\loc(R)$),
并且 $f$ 是典范映射 $X\rightarrow Y=X/R$,
则 $f^{\ast}_r$ 是范畴等价.
因此 $\Top(X/r)$ 是一个 $𝒰$-拓扑斯,
而 $\Top(X/r)\rightarrow\Top(X)$ 是某个拓扑斯态射的逆像函子.

\item 由 i) 得出: 若 $r$ 严格开,
则 $\Top(X/r)$ 是一个拓扑斯,
且包含函子 $\Top(X/r)\rightarrow\Top(X)$
是与某个拓扑斯态射
$$
p:\Top(X)\longrightarrow\Top(X/r){\quoi}.
$$
相应的逆像函子.

\item 对二元组 $(X,r)$,
其中 $X$ 是一个配备严格开局部等价关系的拓扑空间,
定义拓扑斯 $\Top(X/r)$ 以及上述拓扑斯态射 $p$ 的函子性构造.
证明若 $X'\rightarrow X$ 是平展映射,
$r'$ 是 $r$ 在 $X'$ 上的逆像,
则诱导态射 $\Top(X'/r')\rightarrow\Top(X/r)$
等价于局部化态射 $\Top(X/r)_{/E}\rightarrow\Top(X/r)$,
其中 $E$ 是 $\Top(X/r)$ 的一个对象 (在相差唯一同构的意义下确定).
对象 $E$ 覆盖 $\Top(X/r)$ 的终对象,
当且仅当 $X'$ 在 $X$ 中的像在 $R=\glob(r)$ 下的饱和为整个 $X$.

\item 由 k) 得出: $\Top(X/r)$ 是一个\emph{刚性展拓扑斯}
(\hyperref[IV.exer9.8.2]{9.8.2} h)).
证明若 $X$ 清醒,
则 $\Top(X/r)$ 的点的同构类所成的集合,
赋予其典范拓扑 (\hyperref[IV.exer7.8]{7.8}) 后,
同胚于商拓扑空间 $X/\glob(r)$;
并且对于来自 $X$ 的点 $x$ 和 $x'$ 的 $\Top(X/r)$ 的任意两个点,
从其中一个到另一个至多有一个态射.
这样的态射存在, 当且仅当 $x'$ 和 $x$
分别模 $\glob(r)$ 等价于某些点 $x'_1$ 和 $x_1$,
使 $x_1$ 是 $x'_1$ 的一个一般化.

\item 研究典范拓扑斯态射 $\Top(X)\rightarrow\Top(X/r)$,
并注意\pageoriginale, 对 $\Top(X/r)$ 的每个对象 $U$,
若 $U$ 上的诱导拓扑斯“是”通常的拓扑空间,
则诱导态射 $\Top(X)_{/p^{\ast}(U)}\rightarrow\Top(X/r)_{/U}$
“是”一个通常拓扑空间的连续映射 $X_U\rightarrow U$.
证明映射 $X_U\rightarrow U$ 的纤维同胚于 f) 中空间 $X^r$ 的连通分支.

\item 设 $X$ 是一个配备叶状结构的可微流形,
\ie 配备 $X$ 的切层的一个对括号运算稳定、局部为直因子的子层 $F$.
在 $X$ 上定义一个与此叶状结构相应的严格开局部等价关系,
并在商拓扑斯 $T=\Top(X/r)$ 上定义一个环,
使它成为\emph{可微展拓扑斯} (\hyperref[IV.exer9.8.2]{9.8.2} j)).
在可微展拓扑斯 $T$ 上的局部自由模
(也称 $T$ 上的可微向量丛) 所成的范畴,
与 $X$ 上配备相对于给定切丛子丛 $F$ 的联络的局部自由模所成的范畴之间,
定义一个等价.
给出实解析和复解析情形中的变式.
\end{enumerate}
\end{enonce*}

% END INCLUDED SGA4_English_current_buildable_source_20260906/source/04/drafts/SGA4_IV_LINES_8026_8245_English.texfrag


% BEGIN INCLUDED SGA4_English_current_buildable_source_20260906/source/04/drafts/SGA4_IV_LINES_8246_8573_English.texfrag
